% Fragmento de inserción; se incluye desde la raíz de este paquete editorial.
% Las fuentes completas y los cortes materiales constan en ANTECEDENTES_Y_DEPENDENCIAS.md.
% Procedencia: XI REV11 ES, sections/nucleo.tex, tpk_desarrollo_integrado.tex,
% ch_tres_vueltas_ext.tex, ch_elevacion_catalogal.tex, ch_cierre_cardinal.tex,
% ch_descenso_por_fibras.tex; perfil integral public_final/feigenbaum/c37_feigenbaum.tex;
% Grassmannianos REV07 ES, sections/jacobi_spectral.tex;
% TRANSVERSALIDAD_Y_ESCALADO_LOCAL_FEIGENBAUM.md, sección 14.
% Ampliación expositiva de demostraciones existentes; no introduce resultados nuevos.
\providecommand{\ZZ}{\mathbb Z}
\providecommand{\RR}{\mathbb R}
\providecommand{\FF}{\mathbb F}
\providecommand{\APP}{\mathrm{APP}}
\providecommand{\TPK}{\mathrm{TPK}}
\providecommand{\TRIT}{\{-1,0,+1\}}

\section{Antecedentes operatorios de la familia de criticidad}
\label{hmtfg:antecedentes-operatorios}

La familia de criticidad se construye a partir de una publicación del
calendario TPK. Las dos hojas aritméticas, la orientación trítica y el
transporte con memoria determinan primero los estados y sus prolongaciones.
La lectura ordenada del continuo permite después evaluar funciones
analíticas, comparar parámetros y seleccionar raíces. En este desarrollo
se conservan el estado genealógico y sus dos publicaciones: la carta
arquimediana proporciona el espacio de evaluación, mientras la publicación
de incidencia conserva las relaciones de frontera. La estructura discreta
conjunta del continuo mantiene sus cinco construcciones consustanciales.
El presente corte expositivo utiliza su lector ordenado y conserva el
antecedente común del que procede.

\subsection{Células aritméticas, orientación y transporte}
\label{hmtfg:app-trit-transporte}

Se utilizan la célula levantada \eqref{eq:app-levantada}, las coordenadas tríticas y su ley aritmética (\ref{trit-ampl-def-coordenadas}, \ref{trit-ampl-prop-aritmetica}), y el selector~\ref{tpk-int:seleccion}. Para fijar la notación, sean $D_9=\{1,\ldots,9\}$, $X_9=(\mathbb Z/9\mathbb Z)^2$ y, para todo $n\in\mathbb Z$, $\rho_9(n)=1+((n-1)\bmod9)$ y $q_9(n)=(n-\rho_9(n))/9$. El representante balanceado es $r_3(n)\in\{-1,0,+1\}$; su acarreo es $\chi(a,b)=(a+b-r_3(a+b))/3$. Estas fórmulas conservan la extensión a enteros de las evaluaciones positivas del núcleo.

El transporte levantado y su reversibilidad están demostrados en \ref{tpk-int:prop-desplazamiento}; la configuración observable y su retorno de $108$ pasos, en \ref{tpk-int:fobs} y \ref{tpk-int:periodo}. La actualización \eqref{eq:actualizacion} conserva el estado completo. La holonomía de \ref{tpk-int:gamma} y los retornos de \ref{hist:tres-retornos} distinguen restauración de fase y avance de memoria. Sobre estas operaciones se construye a continuación la realización concreta de las tres ramas y su registro íntegro.

% Procedencia literal: output/PREPARACION_ENTREGA_USB_20260928/ENTREGA_HMT_USB/03_FUENTES_Y_REPRODUCCION/ACTUALIZACIONES_20260928/FINAL01/PAQUETES/ARTICULO_XI_REV11_ES_EN_FINAL_BN_20260928/ARTICULO_XI_REV11_ES_EN_FINAL_BN/ES/source/sections/ch_tres_vueltas_ext.tex
% Recuperación íntegra; única transformación mecánica: prefijo hmtfg: en label/ref/eqref.
% No se modifican operadores, pruebas, tablas ni afirmaciones del propietario.
\subsubsection{Realización arista por arista de las tres prolongaciones nonádicas}
\label{hmtfg:subsubsec:tres-vueltas-ext-ch}

Esta subsección especifica el modelo autónomo generado por la relación de
extensión TPK que interviene en la prueba cardinal. No se infiere la existencia de
una arista por coincidencia de tipos: se construyen sus estados inicial y final,
su relación APP, su transporte de fibra y la actualización unaria del registro.
La construcción utiliza exclusivamente las dos hojas de APP, el selector TRIT
y la composición \(\mathsf{Sel}\)--\(\mathsf{Tra}\)--\(\operatorname{Upd}\) del TPK.

\paragraph{Células completas y tres reducciones balanceadas.}
En la carta \(D_9=\{1,\ldots,9\}\), mantendremos siempre la célula APP
completa
\[
 a_{p,q}:=(p,q;R^\Sigma_{p,q},q^\Sigma_{p,q},
                 R^\Pi_{p,q},q^\Pi_{p,q}),
\]
con
\[
 p+q=R^\Sigma_{p,q}+9q^\Sigma_{p,q},\qquad
 pq=R^\Pi_{p,q}+9q^\Pi_{p,q}.
\]
La hoja activa se registra en una coordenada separada; no se escinde
\(a_{p,q}\) en una supuesta célula aditiva y otra multiplicativa. Para
\(r\in D_9\), sean
\begin{equation}
 a(r):=
 \begin{cases}
  a_{1,9},&r=1,\\
  a_{1,r-1},&2\le r\le9,
 \end{cases}
 \qquad
 \tau(r):=M_{\rm ph}(r),
 \qquad
 m(r):=\tau(r)\in\{-1,0,+1\}.
 \label{hmtfg:eq:celula-completa-fase-ch}
\end{equation}
El residuo aditivo de \(a(r)\) proporciona aquí un representante APP de la
fase; esto no convierte \(\Sigma\) en hoja activa universal. El selector
dinámico activa \(\Sigma\) en \(r\in\{1,4,7\}\), activa \(\Pi\) en
\(r\in\{2,5,8\}\) y conserva ambas hojas, sin desplazar cursor, en
\(r\in\{3,6,9\}\). En los tres casos se transportan conjuntamente
\(R^\Sigma,q^\Sigma,R^\Pi,q^\Pi\). Se tiene
\[
 (m(1),\ldots,m(9))=(1,-1,0,1,-1,0,1,-1,0).
\]
Las tres parejas ordenadas por la base nonádica son
\begin{equation}
 \mathfrak p_{+1}:=(1,4),\qquad
 \mathfrak p_{-1}:=(2,5),\qquad
 \mathfrak p_{0}:=(3,6).
 \label{hmtfg:eq:tres-pares-internos-ch}
\end{equation}
Estas parejas son la sección normalizada que toma las dos primeras
ocurrencias de cada tipo TRIT bajo el desplazamiento \(r\mapsto r+3\); el
tercer bloque de fases \(7,8,9\) conserva el control de cierre de la vuelta.
Sus dos miembros pertenecen a una misma fibra de \(M_{\rm ph}\). La
reducción balanceada definida en estos antecedentes da, sin introducir un coeficiente
objetivo,
\begin{equation}
 \begin{array}{c|c|c|c}
 \varepsilon&m(\mathfrak p_\varepsilon^{(1)})
                  +m(\mathfrak p_\varepsilon^{(2)})
   &r_3\bigl(2\varepsilon\bigr)&
   \chi(\varepsilon,\varepsilon)\\ \hline
 +1&1+1=(-1)+3\cdot1&-1&+1\\
 0&0+0=0+3\cdot0&0&0\\
 -1&(-1)+(-1)=1+3\cdot(-1)&+1&-1
 \end{array}
 \label{hmtfg:eq:tres-carry-derivados-ch}
\end{equation}
Por tanto,
\begin{equation}
 \chi(\varepsilon,\varepsilon)
 =\frac{2\varepsilon-r_3(2\varepsilon)}3=\varepsilon.
 \label{hmtfg:eq:carry-no-parametrico-ch}
\end{equation}
El signo de la rama es así la salida del cociclo balanceado aplicado a dos
emisiones TRIT producidas por APP. Cada pareja proporciona una de las tres
salidas de acarreo. Por claridad, indexaremos después la pareja mediante
\(\varepsilon(\mathfrak p):=
\chi(m(\mathfrak p^{(1)}),m(\mathfrak p^{(2)}))\); el símbolo
\(\varepsilon\) no suministra a \(\operatorname{Upd}^{\rm cal}\) un acarreo libre.

La arquitectura de fase, el levantamiento balanceado y el transporte de
memoria se realizan aquí mediante veintisiete incidencias locales y sus
grafos de extensión. La sección normalizada fija representantes de las
fibras de \(M_{\rm ph}\); la cobertura que se demostrará utiliza la existencia
de esos representantes y permanece invariante al sustituirlos por otra
sección compatible.

\paragraph{Datos tipados de las nueve aristas.}
Sea
\(K_{\rm Carr}:=\langle\kappa_{\rm Carr}\rangle_{\rm Ab}\) el módulo de
acarreo. Las acciones están tipadas por
\[
 H_{\rm Carr}:\mathscr P_{\rm TPK}^{\rm op}
   \longrightarrow\operatorname{End}(K_{\rm Carr}),\qquad
 Y_{\rm Carr}:\mathscr P_{\rm TPK}
   \longrightarrow\operatorname{End}(K_{\rm Carr}).
\]
En este modelo generado se toman las acciones triviales
\(H_{\rm Carr}(e)=Y_{\rm Carr}(e)=\operatorname{id}_{K_{\rm Carr}}\): la
hoja activa y la ruta se conservan como coordenadas propias del registro y no
actúan sobre el sumando libre de acarreo. Sea además
\((\mathsf M_d,\jmath_{d',d})\) el sistema directo graduado de memoria. Para
no confundir dos representantes identificados por el colímite, el estado
conserva además el grado como coordenada y se usa el portador etiquetado
\[
 \widetilde{\mathsf M}:=
 \bigsqcup_{d\ge1}\bigl(\{d\}\times\mathsf M_d\bigr),
 \qquad
 \operatorname{gr}_{\mathsf M}(d,\mu)=d.
\]
La publicación en la memoria límite viene dada por
\[
\begin{aligned}
 \mathsf M&:=\varinjlim_d\mathsf M_d,&
 \jmath_d&:\mathsf M_d\longrightarrow\mathsf M,\\
 u&:=\jmath_d(u_d),&
 \iota_{\mathsf M}:\mathbb Z&\longrightarrow\mathsf M,
 \qquad n\longmapsto nu .
\end{aligned}
\]
Aquí \(\jmath_d\) es el mapa canónico del colímite y los \(u_d\) forman una
familia compatible. Su componente libre permite
el lector \(\deg_u(nu)=n\).
Las proyecciones de configuración, fase, memoria, acarreo, ruta y
supervivencia se tiparán sobre el portador calibrado construido
recursivamente más abajo; no se importará para ellas ningún espacio ambiental.
Para cada arista construida \(e\), la
componente calibrada
\(\operatorname{Upd}^{\rm cal}_e:=\mathsf{Mem}^{\rm cal}_e\) denotará la
tercera aplicación de la factorización; su ley se fija campo por campo más
abajo y no se presupone como restricción de una relación ambiental.

Fijemos la profundidad inicial \(k_0=1\), una configuración completa
\(q_\star\in\mathcal Q_{\rm obs}\) de fase uno y
\[
 q_{n,r}:=F_{\rm obs}^{9n+r-1}(q_\star),\qquad
 q_{n,10}:=q_{n+1,1}.
\]
Como el calendario observable es periódico, su valor \(q_{n,r}\) no
distingue por sí solo apariciones separadas por doce vueltas. El modelo
calibrado emplea, por ello, la elevación etiquetada
\[
 \widetilde q_{n,r}:=(q_{n,r},n,r),\qquad
 \widetilde q_{n,10}:=\widetilde q_{n+1,1},\qquad
 \pi_{\rm obs}(\widetilde q_{n,r})=q_{n,r}.
\]
La dinámica elevada
\(\widetilde F_{\rm obs}(\widetilde q_{n,r})
=\widetilde q_{n,r+1}\) proyecta a \(F_{\rm obs}\); la etiqueta no modifica
la fase física, sino que hace distintos sus niveles de aparición.
Se etiqueta la copia de la célula en cada nivel mediante
\[
 a_{n,r}:=(n,r;a(r)),\qquad
 (n,r)^+:=
 \begin{cases}
  (n,r+1),&1\le r<9,\\
  (n+1,1),&r=9.
 \end{cases}
\]
La etiqueta sólo evita identificar dos apariciones de una misma célula; las
dos divisiones euclídeas y sus cuatro coordenadas son las de \(a(r)\).
Para la vuelta \(n\ge0\) y la fase \(r\in D_9\), el símbolo
\[
 e_{\varepsilon,n,r}:\widetilde q_{n,r}
 \longrightarrow\widetilde q_{n,r+1},
 \qquad r^+:=1+(r\bmod9),
\]
ocupa el nivel
\(d_{n,r}:=k_n+r-1\), donde \(k_n:=k_0+9n\). Su registro local contiene
la célula completa \(a(r)\), la hoja activa determinada por \(M_{\rm ph}(r)\),
el tipo \(\tau(r)\),
la marca de pertenencia a \(\mathfrak p_\varepsilon\) y, al alcanzar el segundo miembro
de esa pareja, la identidad certificada
\eqref{hmtfg:eq:tres-carry-derivados-ch}. Definimos sus operadores únicamente a
partir de ese registro:
\begin{align}
 \mathsf C_{\mathsf M}(e_{\varepsilon,n,r})
   &=\jmath_{d_{n,r}+1,d_{n,r}},&
 \kappa_{\mathsf M}(e_{\varepsilon,n,r})
   &=\mathbf1_{\{r=9\}}u_{d_{n,r}+1},
 \label{hmtfg:eq:memoria-arista-ch}\\
       \operatorname{Carr}(e_{\varepsilon,n,r})
   &=d_{\varepsilon,r}\kappa_{\rm Carr},&
 d_{\varepsilon,r}
   &:=\mathbf1_{\{r=\mathfrak p_\varepsilon^{(2)}\}}
       \chi\!\left(m(\mathfrak p_\varepsilon^{(1)}),
                    m(\mathfrak p_\varepsilon^{(2)})\right).
 \label{hmtfg:eq:carry-arista-derivado-ch}
\end{align}
Para las identidades y para toda composición admisible se define, además,
\begin{align}
 \mathsf C_{\mathsf M}(\operatorname{id}_{\widetilde q};d)
   &:=\operatorname{id}_{\mathsf M_d},&
 \kappa_{\mathsf M}(\operatorname{id}_{\widetilde q};d)&:=0,
 \label{hmtfg:eq:memoria-identidad-ch}\\
 \mathsf C_{\mathsf M}(e_t\cdots e_1)
   &:=\mathsf C_{\mathsf M}(e_t)\circ\cdots\circ
      \mathsf C_{\mathsf M}(e_1),&
 \kappa_{\mathsf M}(e_2e_1)
   &:=\mathsf C_{\mathsf M}(e_2)
        \bigl(\kappa_{\mathsf M}(e_1)\bigr)
      +\kappa_{\mathsf M}(e_2).
 \label{hmtfg:eq:cociclo-memoria-calibrado-ch}
\end{align}
La identidad lleva el grado \(d\) porque una misma configuración observable
puede reaparecer en profundidades distintas. La última igualdad, iterada, define
\(\kappa_{\mathsf M}(e_t\cdots e_1)\) para cualquier longitud y es la ley de
cociclo del sistema directo; las inclusiones satisfacen
\(\jmath_{d+2,d+1}\jmath_{d+1,d}=\jmath_{d+2,d}\).
Además se fija
\(s(e_{\varepsilon,n,r})=\top\). Las siguientes veintisiete entradas
certifican las coordenadas que distinguen las tres prolongaciones; los campos
universales del transporte se especifican inmediatamente después. Los símbolos \(I\) y \(II\)
designan, respectivamente, la primera emisión conservada y la segunda emisión
que efectúa la reducción; un guion indica incremento de acarreo nulo, pero no
ausencia de arista.
\begin{center}
\small
\setlength{\tabcolsep}{3.2pt}
\begin{tabular}{c c cc cc cc c}
\toprule
\(r\)&\(M_{\rm ph}(r)\)&\multicolumn{2}{c}{\(\mathfrak p_{+1}\)}&
\multicolumn{2}{c}{\(\mathfrak p_0\)}&
\multicolumn{2}{c}{\(\mathfrak p_{-1}\)}&
\(\Delta\operatorname{mem}\)\\
& &marca&\(d_{+1,r}\)&marca&\(d_{0,r}\)&marca&\(d_{-1,r}\)&\\
\midrule
1&\(+1\)&I&0&--&0&--&0&0\\
2&\(-1\)&--&0&--&0&I&0&0\\
3&0&--&0&I&0&--&0&0\\
4&\(+1\)&II&\(+1\)&--&0&--&0&0\\
5&\(-1\)&--&0&--&0&II&\(-1\)&0\\
6&0&--&0&II&0&--&0&0\\
7&\(+1\)&--&0&--&0&--&0&0\\
8&\(-1\)&--&0&--&0&--&0&0\\
9&0&--&0&--&0&--&0&\(u\)\\
\bottomrule
\end{tabular}
\end{center}
Así, \(d_{\varepsilon,r}\) se calcula a partir de dos salidas del selector y
es distinto de los cocientes \(q^\Sigma,q^\Pi\) de la célula. El registro
conserva también el par marcado y el residuo balanceado; en particular, la
rama nula no se confunde con una ausencia de evento.

La ley de composición del acarreo empleada por la categoría TPK es
\begin{equation}
 \operatorname{Carr}(\operatorname{id}_{\widetilde q})=0,\qquad
 \operatorname{Carr}(e_2e_1)
 =H_{\rm Carr}(e_1)\bigl(\operatorname{Carr}(e_2)\bigr)
  +Y_{\rm Carr}(e_2)\bigl(\operatorname{Carr}(e_1)\bigr).
 \label{hmtfg:eq:ley-carry-correcta-ch}
\end{equation}
En identidades y caminos se prolongan functorialmente:
\[
\begin{aligned}
 H_{\rm Carr}(\operatorname{id})
   &=Y_{\rm Carr}(\operatorname{id})=\operatorname{id},\\
 H_{\rm Carr}(e_t\cdots e_1)
   &=H_{\rm Carr}(e_1)\circ\cdots\circ H_{\rm Carr}(e_t),\\
 Y_{\rm Carr}(e_t\cdots e_1)
   &=Y_{\rm Carr}(e_t)\circ\cdots\circ Y_{\rm Carr}(e_1).
\end{aligned}
\]
En este modelo todas estas acciones siguen siendo la identidad, pero
su tipado explicita qué endomorfismo actúa sobre cada sumando de
\eqref{hmtfg:eq:ley-carry-correcta-ch}.
En las aristas calibradas, \(H_{\rm Carr}=Y_{\rm Carr}=\operatorname{id}\),
por lo que la composición suma los incrementos ya producidos. La ecuación
\eqref{hmtfg:eq:ley-carry-correcta-ch}, y no una identificación entre orientación y
cociente, gobierna el registro acumulado.

\paragraph{Construcción independiente de las tuplas de transición.}
Un estado de la fibra se escribirá, en el orden de sus campos pertinentes,
\[
 x=(a;q;\tau,o,h;b_{\rm ref};
       R^\Sigma,q^\Sigma,R^\Pi,q^\Pi;c;
       \operatorname{Sig};C,\partial C;p_{\APP},p_{\TRIT};
       \lambda;\operatorname{gr}_{\mathsf M};\operatorname{mem};
       s;\operatorname{Led}).
\]
La notación de actualización \(x[f\leftarrow v]\) cambia sólo el campo
indicado. Para \(h\in\{(\Sigma,\Pi),(\Pi,\Sigma)\}\), defínase
\[
 h_r(h)=
 \begin{cases}
  (\Sigma,\Pi),&r\in\{1,4,7\},\\
  (\Pi,\Sigma),&r\in\{2,5,8\},\\
 h,&r\in\{3,6,9\}.
 \end{cases}
\]
La coordenada completa de hoja y aritmética se abrevia por
\[
 \widehat h(x):=
 \bigl(h(x),R^\Sigma(x),q^\Sigma(x),R^\Pi(x),q^\Pi(x)\bigr).
\]
Sea \(\widehat{\mathcal H}^{\rm cal}_{\widetilde q_{n,r}}\) el conjunto de
las quíntuplas
\[
 \bigl(h,R^\Sigma_{a_{n,r}},q^\Sigma_{a_{n,r}},
          R^\Pi_{a_{n,r}},q^\Pi_{a_{n,r}}\bigr),
 \qquad
 h\in\{(\Sigma,\Pi),(\Pi,\Sigma)\},
\]
cuyos cuatro datos aritméticos son las dos divisiones euclídeas de la célula
APP etiquetada \(a_{n,r}\). Esta definición directa no presupone un estado
en un portador todavía no construido. Sean
\(\mathsf{Hist}^{\rm cal}_{\widetilde q}\) el conjunto de caminos finitos
admisibles del grafo de aristas etiquetadas que terminan en
\(\widetilde q\), incluida la identidad, y
\(\mathsf B^{\rm cal}_{\widetilde q}\) el conjunto de sus sucesiones
ordenadas de extremos. Definimos
\[
 \mathsf S^{\rm cal}_{\widetilde q}:=
 \widehat{\mathcal H}^{\rm cal}_{\widetilde q}
 \times\{o_{\rightarrow},o_{\circ},o_{\leftarrow}\}
 \times K_{\rm Carr}\times
 \mathsf{Hist}^{\rm cal}_{\widetilde q}\times
 \mathsf B^{\rm cal}_{\widetilde q}\times\mathsf M .
\]
La firma calibrada no se postula: es la aplicación de reunión tipada
\[\begin{gathered}
 \operatorname{Sig}^{\rm cal}_{\widetilde q}:
 \widehat{\mathcal H}^{\rm cal}_{\widetilde q}
 \times\{o_{\rightarrow},o_{\circ},o_{\leftarrow}\}
 \times K_{\rm Carr}\times
 \mathsf{Hist}^{\rm cal}_{\widetilde q}\times
 \mathsf B^{\rm cal}_{\widetilde q}\times\mathsf M
 \longrightarrow\mathsf S^{\rm cal}_{\widetilde q},
\\
 (\widehat h,o,c,\lambda,\partial C,\mu)
 \longmapsto(\widehat h,o,c,\lambda,\partial C,\mu).
\end{gathered}\]
Asimismo, \(\pi_{\mathcal F}\) denota la proyección que elimina únicamente
las dos primeras coordenadas \((a;q)\) de la tupla de estado; conserva
\(\tau,o,h,b_{\rm ref}\), los cuatro datos euclídeos, acarreo, firma,
cilindro, frontera, prefijos, ruta, grado de memoria, memoria, supervivencia y
registro.
La orientación correspondiente queda fijada por
\(o(+1)=o_{\rightarrow}\), \(o(0)=o_{\circ}\) y
\(o(-1)=o_{\leftarrow}\).
La fibra APP calibrada sobre \(\widetilde q_{n,r}\) contendrá la copia etiquetada
\(a_{n,r}\) ya definida. El estado inicial queda fijado sin coordenadas libres:
\[
 \widehat h_\star:=
 \bigl((\Sigma,\Pi),R^\Sigma_{1,9},q^\Sigma_{1,9},
       R^\Pi_{1,9},q^\Pi_{1,9}\bigr).
\]
\begin{equation}
\begin{aligned}
 x_\star:=\bigl(&a_{0,1};\widetilde q_{0,1};
 +1,o_{\rightarrow},(\Sigma,\Pi);\bot;
 R^\Sigma_{1,9},q^\Sigma_{1,9},R^\Pi_{1,9},q^\Pi_{1,9};0;\\
 &\operatorname{Sig}^{\rm cal}_{\widetilde q_{0,1}}
       (\widehat h_\star,o_{\rightarrow},0,
        \operatorname{id}_{\widetilde q_{0,1}},\varnothing,
        \jmath_{k_0}(0_{\mathsf M_{k_0}}));
 \varnothing,\varnothing;
 (a_{0,1}),(+1);\operatorname{id}_{\widetilde q_{0,1}};
 k_0;0_{\mathsf M_{k_0}};\top;\varnothing\bigr).
\end{aligned}
\label{hmtfg:eq:estado-inicial-calibrado-ch}
\end{equation}
La tupla no contiene coordenadas indeterminadas y satisface
\(s(x_\star)=\top\). Sea \(G_0:=\{x_\star\}\).
Supuesto construido \(G_n\), fijemos \(x\in G_n\),
\(\varepsilon\in\{-1,0,+1\}\) y
\(x_{\varepsilon,0}(x):=x\). Para \(r=1,\ldots,9\), se construye primero
el registro elemental
\begin{equation}
\begin{aligned}
 \ell_{\varepsilon,n,r}:=\bigl(&e_{\varepsilon,n,r},a_{n,r},a_{(n,r)^+},
 \widetilde q_{n,r},\widetilde q_{n,r+1},M_{\rm ph}(r),\\
 &h_r(h(x_{\varepsilon,r-1}(x))),o(M_{\rm ph}(r)),\mathfrak p_\varepsilon,
 d_{\varepsilon,r},\\
 &\mathsf C_{\mathsf M}(e_{\varepsilon,n,r}),
 \kappa_{\mathsf M}(e_{\varepsilon,n,r}),
 \operatorname{Carr}(e_{\varepsilon,n,r})\bigr),
\end{aligned}
 \label{hmtfg:eq:registro-elemental-ext-ch}
\end{equation}
y se pone
\(\operatorname{Led}^{\rm new}_{\varepsilon,n,r}:=
\operatorname{Led}(x_{\varepsilon,r-1}(x))^\frown
\ell_{\varepsilon,n,r}\). Se construye entonces la tupla
\begin{equation}
 y_{\varepsilon,r}(x):=
 x_{\varepsilon,r-1}(x)
 [\tau\leftarrow M_{\rm ph}(r),\
  h\leftarrow h_r(h(x_{\varepsilon,r-1}(x))),\
  o\leftarrow o(M_{\rm ph}(r)),\
  b_{\rm ref}\leftarrow\mathfrak p_\varepsilon],
 \label{hmtfg:eq:tupla-sel-cal-ch}
\end{equation}
donde la selección deja intactos célula, cocientes, acarreo, ruta, memoria,
cilindro, frontera, supervivencia y registro; la coordenada
\(b_{\rm ref}\) hace explícita la reducción balanceada escogida. No se presupone aquí
pertenencia a una relación ambiental.

A continuación se define \(z_{\varepsilon,r}(x)\) campo por campo: su
configuración es \(\widetilde q_{n,r+1}\); su célula completa es
\(a_{(n,r)^+}\), con los
cuatro residuos y cocientes recalculados por las dos divisiones euclídeas;
conserva la hoja seleccionada y la orientación; y efectúa las seis
actualizaciones siguientes, en las que sólo se omite el argumento común \(x\):
\begin{equation}
\begin{aligned}
 C(z_{\varepsilon,r})&:=C(y_{\varepsilon,r})^\frown
       e_{\varepsilon,n,r},\\
 \partial C(z_{\varepsilon,r})&:=\partial C(y_{\varepsilon,r})^\frown
       (\widetilde q_{n,r},\widetilde q_{n,r+1}),\\
 p_{\APP}(z_{\varepsilon,r})&:=p_{\APP}(y_{\varepsilon,r})^\frown
       a_{(n,r)^+},\\
 p_{\TRIT}(z_{\varepsilon,r})&:=p_{\TRIT}(y_{\varepsilon,r})^\frown
       \tau(r),\\
 \operatorname{Led}(z_{\varepsilon,r})&:=
       \operatorname{Led}^{\rm new}_{\varepsilon,n,r},\\
 \lambda(z_{\varepsilon,r})&:=e_{\varepsilon,n,r}
       \lambda(y_{\varepsilon,r}),
\end{aligned}
 \label{hmtfg:eq:tupla-tra-cal-ch}
\end{equation}
y mantiene todavía los campos
\(c,\operatorname{gr}_{\mathsf M},\operatorname{mem},s\) de
\(y_{\varepsilon,r}(x)\). Su firma, en cambio, se vuelve a tipar en la
configuración de destino mediante la prefirma transportada
\begin{equation}
\begin{aligned}
 \operatorname{Sig}(z_{\varepsilon,r}(x))
  :=\operatorname{Sig}^{\rm cal}_{\widetilde q_{n,r+1}}
       \Bigl(&\widehat h(z_{\varepsilon,r}(x)),
             o(z_{\varepsilon,r}(x)),c(y_{\varepsilon,r}(x)),
             \lambda(z_{\varepsilon,r}(x)),\\[-2pt]
             &\partial C(z_{\varepsilon,r}(x)),
             \jmath_{d_{n,r}}
                (\operatorname{mem}_{d_{n,r}}
                   (y_{\varepsilon,r}(x)))\Bigr).
\end{aligned}
\label{hmtfg:eq:prefirma-transporte-cal-ch}
\end{equation}
Así, y sólo en
\(\mathsf{Tra}^{\rm cal}\), se añaden la arista, sus extremos, la marca
\((\mathfrak p_\varepsilon,r,d_{\varepsilon,r})\) y sus cociclos al registro de
incidencias.

El estado inicial satisface
\[
 c(x_\star)=0=\operatorname{Carr}
   (\operatorname{id}_{\widetilde q_{0,1}})
   =\operatorname{Carr}(\lambda(x_\star)).
\]
En los estados completos \(x_{\varepsilon,r-1}(x)\), y también después de la
selección \(y_{\varepsilon,r}(x)\), se conserva el invariante tipado
\(c(v)=\operatorname{Carr}(\lambda(v))\), pues
\(\mathsf{Sel}^{\rm cal}\) no modifica ninguno de esos dos campos. El
transporte avanza la ruta pero deja pendiente el acarreo; por ello la tupla de
preactualización satisface exactamente
\[
 c(z_{\varepsilon,r}(x))=c(y_{\varepsilon,r}(x))
 =\operatorname{Carr}(\lambda(y_{\varepsilon,r}(x))).
\]
La aplicación \(\operatorname{Upd}^{\rm cal}\) restaura entonces
\(c(x_{\varepsilon,r}(x))=
\operatorname{Carr}(\lambda(x_{\varepsilon,r}(x)))\) mediante la ley de
composición, sin atribuir ese invariante al estado intermedio \(z\).

Finalmente se aplica la función unaria
\(\operatorname{Upd}^{\rm cal}_{e_{\varepsilon,n,r}}
=\mathsf{Mem}^{\rm cal}_{e_{\varepsilon,n,r}}\) a esa tupla. Su
resultado \(x_{\varepsilon,r}(x)\) viene dado por
\begin{equation}
\begin{aligned}
 c\bigl(x_{\varepsilon,r}(x)\bigr)
   &=\operatorname{Carr}\bigl(\lambda(z_{\varepsilon,r}(x))\bigr)\\
   &=H_{\rm Carr}(\lambda(y_{\varepsilon,r}(x)))
       \bigl(\operatorname{Carr}(e_{\varepsilon,n,r})\bigr)
     +Y_{\rm Carr}(e_{\varepsilon,n,r})
       \bigl(c(y_{\varepsilon,r}(x))\bigr),\\
 \lambda\bigl(x_{\varepsilon,r}(x)\bigr)
   &=\lambda(z_{\varepsilon,r}(x)),\\
 \operatorname{gr}_{\mathsf M}
      \bigl(x_{\varepsilon,r}(x)\bigr)
   &=d_{n,r}+1,\\
 \operatorname{mem}_{d_{n,r}+1}
       \bigl(x_{\varepsilon,r}(x)\bigr)
   &=\mathsf C_{\mathsf M}(e_{\varepsilon,n,r})
       \operatorname{mem}_{d_{n,r}}(z_{\varepsilon,r}(x))
     +\kappa_{\mathsf M}(e_{\varepsilon,n,r}),\\
 s\bigl(x_{\varepsilon,r}(x)\bigr)
   &=s(e_{\varepsilon,n,r})\wedge s(z_{\varepsilon,r}(x)),\\
 \operatorname{Sig}\bigl(x_{\varepsilon,r}(x)\bigr)
   &=\operatorname{Sig}^{\rm cal}_{\widetilde q_{n,r+1}}
       \Bigl(\widehat h(z_{\varepsilon,r}(x)),
             o(z_{\varepsilon,r}(x)),
             c(x_{\varepsilon,r}(x)),
             \lambda(z_{\varepsilon,r}(x)),
             \partial C(z_{\varepsilon,r}(x)),\\[-2pt]
   &\hspace{112pt}
             \jmath_{d_{n,r}+1}
              \bigl(\operatorname{mem}_{d_{n,r}+1}
                 (x_{\varepsilon,r}(x))\bigr)\Bigr).
\end{aligned}
\label{hmtfg:eq:upd-unaria-tres-vueltas-ch}
\end{equation}
Todos los campos no enumerados en
\eqref{hmtfg:eq:upd-unaria-tres-vueltas-ch} coinciden con los de
\(z_{\varepsilon,r}(x)\). En particular, la firma se calcula con la hoja,
orientación, frontera y ruta ya fijadas por \(\mathsf{Sel}^{\rm cal}\) y
\(\mathsf{Tra}^{\rm cal}\), y con el acarreo y la memoria calculados en las
líneas anteriores; no se define circularmente a partir del estado final.
Ni \(\varepsilon\), ni \(d_{\varepsilon,r}\), ni \(u\) son argumentos
adicionales de \(\operatorname{Upd}^{\rm cal}\): la función lee los tres datos de la
arista ya construida.

\paragraph{Inducción simultánea de los niveles.}
Completadas las nueve etapas anteriores para un \(G_n\) ya construido,
se define inmediatamente
\begin{equation}
 \gamma_{\varepsilon,k_n}(x):=x_{\varepsilon,9}(x),\qquad
 G_{n+1}:=\{\gamma_{\varepsilon,k_n}(x):
       x\in G_n,\ \varepsilon\in\TRIT\}.
 \label{hmtfg:eq:injerto-total-tres-vueltas-ch}
\end{equation}
\label{hmtfg:eq:tres-vueltas-tpk-ch}
La base \(G_0=\{x_\star\}\), las tuplas intermedias y
\eqref{hmtfg:eq:injerto-total-tres-vueltas-ch} constituyen una única definición por
recursión sobre \(n\). Por tanto, ningún portador posterior se utiliza para
suponer la existencia de \(G_n\): el nivel siguiente sólo se introduce
después de construir todas sus aristas desde el nivel precedente.

\paragraph{Relación de extensión TPK y pertenencia efectiva.}
Sea \(\mathcal E^{\rm cal}_{\widetilde q}\) la menor familia etiquetada que contiene
\(x_\star\) y, en cada paso de la recurrencia anterior, las cuatro tuplas
\(x_{\varepsilon,r-1}(x),y_{\varepsilon,r}(x),
z_{\varepsilon,r}(x),x_{\varepsilon,r}(x)\) en sus configuraciones de origen
y destino. Sea asimismo
\(\mathcal A^{\rm cal}_{\widetilde q_{n,r}}:=\{a_{n,r}\}\), y póngase
\[
 \widetilde{\mathcal Q}_{\rm obs}:=
   \{\widetilde q_{n,r}:n\ge0,\ 1\le r\le9\},
 \qquad
 \mathcal E^{\rm cal}:=\bigsqcup_{\widetilde q}
     \mathcal E^{\rm cal}_{\widetilde q},
 \qquad
 \mathcal F^{\rm cal}_{\widetilde q}:=
     \pi_{\mathcal F}(\mathcal E^{\rm cal}_{\widetilde q}).
\]
Para cada profundidad \(d\), sea
\(\mathcal E^{\rm cal}_d:=
\operatorname{gr}_{\mathsf M}^{-1}(d)\subset\mathcal E^{\rm cal}\). La
coordenada explícita de grado, y no la mera pertenencia de un representante a
alguna imagen del sistema directo, hace disjuntos estos portadores. Sobre ellos
quedan tipadas las proyecciones
\[
\begin{aligned}
 \operatorname{cfg}&:\mathcal E^{\rm cal}\longrightarrow
     \widetilde{\mathcal Q}_{\rm obs},\\
 g(x)&:=\operatorname{fase}\!
       \bigl(\pi_{\rm obs}(\operatorname{cfg}(x))\bigr)\in C_9,\\
 c&:\mathcal E^{\rm cal}\longrightarrow K_{\rm Carr},\\
 s&:\mathcal E^{\rm cal}\longrightarrow\{\bot,\top\},\\
 \lambda&:\mathcal E^{\rm cal}_{\widetilde q}
       \longrightarrow\mathsf{Hist}^{\rm cal}_{\widetilde q},\\
 \operatorname{gr}_{\mathsf M}&:\mathcal E^{\rm cal}
       \longrightarrow\mathbb N_{\ge1},\\
 \operatorname{mem}_d&:\mathcal E^{\rm cal}_d
       \longrightarrow\mathsf M_d,\\
 \operatorname{mem}(x)&:=
       \jmath_d\bigl(\operatorname{mem}_d(x)\bigr)\in\mathsf M
       \quad(x\text{ de profundidad }d).
\end{aligned}
\]
Aquí \(\operatorname{cfg}(x)\) es la segunda coordenada de la tupla y
\(\operatorname{fase}\) lee la fase del calendario observable. Para las tres
clases de portador se fijan las diagonales
\[
\begin{aligned}
 \Delta_{\APP,\widetilde q}
   &:=\{(a,a):a\in\mathcal A^{\rm cal}_{\widetilde q}\},\\
 \Delta_{{\rm fib},\widetilde q}
   &:=\{(f,f):f\in\mathcal F^{\rm cal}_{\widetilde q}\},\\
 \Delta_{{\rm enr},\widetilde q}
   &:=\{(x,x):x\in\mathcal E^{\rm cal}_{\widetilde q}\}.
\end{aligned}
\]
En particular,
\[
 a_{0,1}\in\mathcal A^{\rm cal}_{\widetilde q_{0,1}},
 \qquad
 x_\star\in\mathcal E^{\rm cal}_{\widetilde q_{0,1}}.
\]
Sobre estos portadores se
definen por sus grafos completos, y no mediante implicaciones necesarias, las
relaciones
\begin{align}
 \mathsf{Sel}^{\rm cal}_{e_{\varepsilon,n,r}}
   &:=\{(x_{\varepsilon,r-1}(x),y_{\varepsilon,r}(x)):
          x\in G_n\},\\
 \mathsf{Tra}^{\rm cal}_{e_{\varepsilon,n,r}}
   &:=\{(y_{\varepsilon,r}(x),z_{\varepsilon,r}(x)):
          x\in G_n\},\\
 \operatorname{Ext}^{\APP,\rm cal}_{e_{\varepsilon,n,r}}
   &:=\{(a_{n,r},a_{(n,r)^+})\},
 \label{hmtfg:eq:ext-app-calibrada-ch}\\
 \operatorname{Ext}^{\rm fib,cal}_{e_{\varepsilon,n,r}}
   &:=\{(\pi_{\mathcal F}x_{\varepsilon,r-1}(x),
          \pi_{\mathcal F}x_{\varepsilon,r}(x)):x\in G_n\}.
 \label{hmtfg:eq:ext-fibra-calibrada-ch}
\end{align}
La extensión enriquecida sincronizada es
\begin{equation}
 \operatorname{Ext}^{\rm enr,cal}_{e_{\varepsilon,n,r}}
 :=\{(x_{\varepsilon,r-1}(x),x_{\varepsilon,r}(x)):
       x\in G_n\}.
 \label{hmtfg:eq:ext-enriquecida-calibrada-ch}
\end{equation}
La tercera aplicación queda definida por
\begin{equation}
 \operatorname{graph}
 \bigl(\operatorname{Upd}^{\rm cal}_{e_{\varepsilon,n,r}}\bigr)
 :=\{(z_{\varepsilon,r}(x),x_{\varepsilon,r}(x)):
       x\in G_n\},
 \label{hmtfg:eq:grafo-upd-calibrado-ch}
\end{equation}
donde el segundo componente es exactamente el de
\eqref{hmtfg:eq:upd-unaria-tres-vueltas-ch}. Por consiguiente,
\begin{equation}
 \mathcal U^{(1),\rm cal}_{e_{\varepsilon,n,r}}
 :=\operatorname{Upd}^{\rm cal}_{e_{\varepsilon,n,r}}\circ
   \mathsf{Tra}^{\rm cal}_{e_{\varepsilon,n,r}}\circ
   \mathsf{Sel}^{\rm cal}_{e_{\varepsilon,n,r}}
 =\{(x_{\varepsilon,r-1}(x),x_{\varepsilon,r}(x)):
      x\in G_n\}.
 \label{hmtfg:eq:transicion-cal-en-tpk-ch}
\end{equation}
Cada operador está indexado por la arista
\(e_{\varepsilon,n,r}\) producida por la reducción
\(\mathfrak p_\varepsilon\), y no por un trit suministrado desde fuera. La
coordenada \(b_{\rm ref}\) distingue además sus codominios. Por ello
\(\mathsf{Sel}^{\rm cal}_{e}\),
\(\mathsf{Tra}^{\rm cal}\) y \(\operatorname{Upd}^{\rm cal}\) son funciones
sobre sus respectivos dominios, mientras
la unión
\(\bigcup_{\varepsilon\in\TRIT}
\mathcal U^{(1),\rm cal}_{e_{\varepsilon,n,r}}\) conserva deliberadamente
las tres salidas de la trisección.

Para \(\bullet\in\{\APP,{\rm fib},{\rm enr}\}\), las identidades y las
composiciones no quedan implícitas. Se definen
\begin{align}
 \operatorname{Ext}^{\bullet,\rm cal}_{\operatorname{id}_{\widetilde q}}
   &:=\Delta_{\bullet,\widetilde q},\\
 \operatorname{Ext}^{\bullet,\rm cal}_{e_t\cdots e_1}
   &:=\operatorname{Ext}^{\bullet,\rm cal}_{e_t}\circ\cdots\circ
      \operatorname{Ext}^{\bullet,\rm cal}_{e_1}
 \label{hmtfg:eq:ext-palabra-calibrada-ch}
\end{align}
para toda palabra admisible de aristas. Las dos clausuras de caminos APP y
fibra se forman por separado,
\begin{equation}
 \operatorname{Path}(\operatorname{Ext}^{\APP,\rm cal}),\qquad
 \operatorname{Path}(\operatorname{Ext}^{\rm fib,cal}),
 \label{hmtfg:eq:dos-clausuras-ext-ch}
\end{equation}
y \(\mathsf{Tra}^{\rm cal}\) las sincroniza arista por arista.

Finalmente, los niveles de estados completos y sus truncamientos se fijan
por la propia recurrencia:
\begin{align}
 H^{\rm cal}_{d_{n,r}}
   &:=\{x_{\varepsilon,r-1}(x):
          x\in G_n,\ \varepsilon\in\TRIT\},\\
 H^{\rm cal}_{d_{n,r}+1}
   &:=\{x_{\varepsilon,r}(x):
          x\in G_n,\ \varepsilon\in\TRIT\},\\
 \rho^{\rm cal}_{d_{n,r}+1,d_{n,r}}
   \bigl(x_{\varepsilon,r}(x)\bigr)
   &:=x_{\varepsilon,r-1}(x).
 \label{hmtfg:eq:truncamiento-calibrado-definido-ch}
\end{align}
El último registro conserva \((n,r,\mathfrak p_\varepsilon)\), de modo que
el predecesor del miembro izquierdo es único y \(\rho^{\rm cal}\) está bien
definida. Para \(b>a\), se pone
\(\rho^{\rm cal}_{b,a}:=\rho^{\rm cal}_{a+1,a}\circ\cdots\circ
\rho^{\rm cal}_{b,b-1}\). Esta aplicación recupera el estado anterior
completo, incluida su firma; no pretende invertir una coordenada sobrescrita
sin consultar el registro.

\begin{proposition}[Realización autónoma del TPK mediante transiciones explícitas]
\label{hmtfg:prop:modelo-tpk-calibrado-ch}
La familia
\[
\begin{aligned}
 \mathfrak T_{\rm TPK}^{\rm cal}:=\bigl(
 &(\widetilde q_{n,r})_{n,r},\pi_{\rm obs},
   \widetilde F_{\rm obs},M_{\rm ph},\\
 &(\mathcal A_{\widetilde q}^{\rm cal})_{\widetilde q},
   (\mathcal E_{\widetilde q}^{\rm cal})_{\widetilde q},
   (\mathcal F_{\widetilde q}^{\rm cal})_{\widetilde q},
   \operatorname{Ext}^{\APP,\rm cal},
   \operatorname{Ext}^{\rm fib,cal},
   \operatorname{Ext}^{\rm enr,cal},\rho^{\rm cal},\\
 &\mathsf{Sel}^{\rm cal},\mathsf{Tra}^{\rm cal},
   \operatorname{Upd}^{\rm cal},\mathsf C_{\mathsf M},
   \kappa_{\mathsf M},\operatorname{gr}_{\mathsf M},
   H_{\rm Carr},Y_{\rm Carr},\\
 &\operatorname{Sig}^{\rm cal},\pi_{\mathcal F}\bigr).
\end{aligned}
\]
es un modelo enriquecido del TPK. En particular, cada par del miembro derecho
de \eqref{hmtfg:eq:transicion-cal-en-tpk-ch} es una arista TPK efectiva, y no la
consecuencia inferida de una condición solamente necesaria.
\end{proposition}

\begin{proof}
Las etiquetas \((n,r)\) hacen disjuntas las copias de los portadores y fijan
sin ambigüedad cada origen y destino. La relación
\(\operatorname{Ext}^{\APP,\rm cal}\) transporta la célula completa y
preserva las dos identidades euclídeas; la relación
\(\operatorname{Ext}^{\rm fib,cal}\) relaciona las fibras de los estados
completos anterior y posterior: transporta hoja, orientación, prefijo,
cilindro, frontera y registro e incorpora el acarreo, la memoria y la firma
actualizados. La selección no modifica el registro y el transporte concatena
exactamente una entrada y produce la prefirma de destino
\eqref{hmtfg:eq:prefirma-transporte-cal-ch}. La actualización recalcula la firma mediante la
aplicación tipada \(\operatorname{Sig}^{\rm cal}_{\widetilde q}\) desde
hoja, orientación, acarreo, ruta, frontera y memoria ya determinados. La aplicación
\(\operatorname{Upd}^{\rm cal}\) es unaria porque la arista contenida en
\(z_{\varepsilon,r}(x)\) determina sus cociclos.

Las identidades son las diagonales y las extensiones por palabras son las
composiciones de \eqref{hmtfg:eq:ext-palabra-calibrada-ch}. La concatenación de
prefijos y registros es asociativa; la memoria satisface la ley de cociclo y
el acarreo satisface
\eqref{hmtfg:eq:ley-carry-correcta-ch}. Por tanto, las dos parentizaciones de tres
aristas producen el mismo estado. La fórmula
\eqref{hmtfg:eq:truncamiento-calibrado-definido-ch} define el truncamiento sobre
los estados completos y la unicidad del último registro prueba que recupera
el origen en todos los campos. Éstas son las leyes de identidad, fuente,
destino, composición, memoria, acarreo y truncamiento del TPK enriquecido.
En particular, las relaciones de extensión quedan determinadas como grafos
de aplicaciones efectivas y no como consecuencias de meras condiciones
necesarias.
\end{proof}

Cada \(\gamma_{\varepsilon,k_n}\) comprende exactamente nueve aplicaciones
de \(\operatorname{Upd}^{\rm cal}_{e_{\varepsilon,n,r}}\circ
\mathsf{Tra}^{\rm cal}_{e_{\varepsilon,n,r}}\circ
\mathsf{Sel}^{\rm cal}_{e_{\varepsilon,n,r}}\). La
supervivencia se probará mediante una cola construida y no se utilizará para
definir \(G_n\).

\begin{lemma}[Trisección interna efectiva de una vuelta]
\label{hmtfg:lem:elevacion-catalogal-nonadica}
Para todo \(n\ge0\), todo \(x\in G_n\) y todo
\(\varepsilon\in\{-1,0,+1\}\), la expresión
\(\gamma_{\varepsilon,k_n}(x)\) está definida por nueve transiciones TPK,
con sus relaciones sincronizadas
\(\operatorname{Ext}^{\APP,\rm cal}\) y
\(\operatorname{Ext}^{\rm fib,cal}\), y satisface
\begin{align}
 \rho^{\rm cal}_{k_n+9,k_n}\gamma_{\varepsilon,k_n}(x)&=x,
 \label{hmtfg:eq:truncamiento-injerto-ch}\\
 g\bigl(\gamma_{\varepsilon,k_n}(x)\bigr)
   &=g(x),&
 \jmath_{k_n+9}\!\left(\operatorname{mem}_{k_n+9}
        (\gamma_{\varepsilon,k_n}(x))\right)
   &=\jmath_{k_n}\!\left(\operatorname{mem}_{k_n}(x)\right)+u,
 \label{hmtfg:eq:fase-memoria-injerto-ch}\\
 c\bigl(\gamma_{\varepsilon,k_n}(x)\bigr)
   &=c(x)+\varepsilon\kappa_{\rm Carr}.&&
 \label{hmtfg:eq:carry-injerto-ch}
\end{align}
\label{hmtfg:eq:tres-vueltas-propiedades-ch}
\label{hmtfg:eq:tres-vueltas-acarreo-ch}
El registro ordenado de las nueve aristas determina \(\varepsilon\) de manera única.
Las tres imágenes son, por tanto, disjuntas. Cada extremo pertenece a
\(G_{n+1}\) y vuelve a admitir las tres aplicaciones
\(\gamma_{-1,k_{n+1}},\gamma_{0,k_{n+1}},
\gamma_{+1,k_{n+1}}\).
\end{lemma}

\begin{proof}
Las fórmulas \eqref{hmtfg:eq:celula-completa-fase-ch} recorren una vez las nueve
fases y devuelven la copia \(a_{n+1,1}\) de \(a(1)\). En cada paso, la relación APP transporta el
séxtuplo completo; \(\mathsf{Sel}^{\rm cal}\) fija el tipo TRIT desde la marca de fase
\(r\) conservada en el estado;
y \(\mathsf{Tra}^{\rm cal}\) añade exactamente un símbolo al prefijo, un tramo al
cilindro y su par ordenado de extremos a la frontera. Por inducción sobre
\(r\), las tuplas explícitas \(y_{\varepsilon,r}(x)\) y
\(z_{\varepsilon,r}(x)\) satisfacen los predicados de las relaciones
indicadas, y la actualización unaria produce el estado del nivel siguiente.
La identidad
\(\rho^{\rm cal}_{d_{n,r}+1,d_{n,r}}
(x_{\varepsilon,r}(x))=x_{\varepsilon,r-1}(x)\) es la definición
\eqref{hmtfg:eq:truncamiento-calibrado-definido-ch}; su buena definición procede de
la unicidad del último registro. La composición de las nueve aplicaciones de
truncamiento recupera \(x\), lo que prueba
\eqref{hmtfg:eq:truncamiento-injerto-ch} sin invertir por separado los campos
sobrescritos.

La fase \(9\to1\) aparece una sola vez. Por
\eqref{hmtfg:eq:memoria-arista-ch} y la ley de cociclo, la parte lineal de memoria
transporta la memoria previa y su parte traslacional suma exactamente \(u\)
en el límite directo. La compatibilidad
\(\jmath_{d+1}\jmath_{d+1,d}=\jmath_d\) prueba
\eqref{hmtfg:eq:fase-memoria-injerto-ch}; el retorno concierne a la fase, no al
estado enriquecido.

Por \eqref{hmtfg:eq:carry-arista-derivado-ch}, ocho aristas poseen incremento nulo y
la arista que completa \(\mathfrak p_\varepsilon\) posee el incremento calculado
\(\chi(\varepsilon,\varepsilon)\kappa_{\rm Carr}\). La ley
\eqref{hmtfg:eq:ley-carry-correcta-ch} y
\eqref{hmtfg:eq:carry-no-parametrico-ch} dan
\eqref{hmtfg:eq:carry-injerto-ch}. El registro conserva
\((\mathfrak p_\varepsilon,2\varepsilon,r_3(2\varepsilon),
\chi(\varepsilon,\varepsilon))\); los tres valores de este registro son
distintos incluso cuando el incremento escalar es cero. Por ello las imágenes
no se identifican y el decodificador de bloque recupera una única
\(\varepsilon\).

Las definiciones de las aristas dependen sólo de la nueva fase y concatenan,
en vez de sustituir, prefijo, cilindro, frontera y registro. También son
invariantes bajo
\(\jmath_d\operatorname{mem}_d\mapsto
  \jmath_d\operatorname{mem}_d+nu\)
y bajo la traslación del acarreo de
entrada. En consecuencia, la misma construcción se aplica al extremo de
cualquier vuelta. Para cada estado finito, la sucesión explícita
\[
 x,\quad\gamma_{0,k_n}(x),\quad
 \gamma_{0,k_{n+1}}\gamma_{0,k_n}(x),\quad\ldots
\]
es una cola infinita compatible. Esto demuestra la supervivencia de todos los
estados construidos y la disponibilidad reiterada de las tres ramas, sin
postular ninguna de ambas propiedades. Por consiguiente podemos, finalmente,
identificar, de acuerdo con las definiciones de nivel anteriores,
\begin{equation}
 H_{k_n}^{\rm cal}:=G_n
 \label{hmtfg:eq:subarbol-calibrado-superviviente-ch}
\end{equation}
para todo \(n\ge0\).
\end{proof}

El cierre anterior distingue los tres relojes implicados. Nueve emisiones
devuelven la fase en \(C_9\) y avanzan la memoria. Doce vueltas, es decir,
\(108\) emisiones, devuelven además la posición del calendario observable;
tampoco entonces se borran el registro ni la memoria. El operador
\(\mathcal K_{\rm ph}\) actúa únicamente después, como reconocimiento de
memoria reducida de las vueltas ya construidas; nunca genera las aristas de
\(\operatorname{Ext}^{\rm enr,cal}\).


\subsection{Refinamiento de cilindros y lectura ordenada}
\label{hmtfg:refinamiento-ordenado}

% XI REV11, ch_elevacion_catalogal.tex: bloques de 54 aristas, beta_blk,
% partición efectiva de 729 subcilindros; ch_cierre_cardinal.tex:375--497.
Para \(\mathbf b=(b_1,\ldots,b_6)\in\FF_3^6\), sea
\(\upsilon:\FF_3\to\{-1,0,+1\}\) el representante balanceado.
Las tres prolongaciones construidas arista por arista definen
\[
 \Gamma^{[54]}_{\mathbf b,k}
 =\gamma_{\upsilon(b_6),k+45}\circ\cdots\circ
   \gamma_{\upsilon(b_1),k}.
\]
Partiendo de \(Y_0=\{x_\star\}\), se pone
\[
 Y_{N+1}
 =\{\Gamma^{[54]}_{\mathbf b,1+54N}(x):
          x\in Y_N,\ \mathbf b\in\FF_3^6\}.
\]
El truncamiento \(\rho^{\rm blk}_{N+1,N}\) elimina las últimas
cincuenta y cuatro aristas, recuperando el estado anterior mediante
el registro de transición.

\begin{proposition}[Codificación efectiva de los bloques de prolongación]
\label{hmtfg:bloques-729}
La lectura de las seis parejas tríticas de cada bloque define biyecciones
\[
 \beta_N:Y_N\xrightarrow{\ \cong\ }(\FF_3^6)^N
\]
que conmutan con el truncamiento. Cada estado de \(Y_N\) posee
exactamente \(729\) prolongaciones de bloque en \(Y_{N+1}\).
\end{proposition}
\begin{proof}
En cada vuelta, el registro conserva la pareja trítica y su acarreo.
Las tres parejas producen las tres salidas distintas
\(\chi(\varepsilon,\varepsilon)=\varepsilon\), y el registro permite
recuperar la rama utilizada, incluida la rama de acarreo nulo.
Las seis concatenaciones producen, por tanto, todas las palabras de
\(\FF_3^6\). Dos palabras diferentes se distinguen en la primera pareja
registrada en la que difieren. La recuperación del prefijo prueba la
compatibilidad con el truncamiento. Por inducción sobre \(N\), cada
palabra de longitud \(N\) tiene un único estado antecedente; la siguiente
palabra puede ser cualquiera de las \(3^6=729\) posibilidades.
\end{proof}

Sea \(\Omega_{\Gamma}^{\rm cal}=\varprojlim_NY_N\). Las biyecciones
anteriores inducen
\[
 \Omega_{\Gamma}^{\rm cal}\cong(\FF_3^6)^{\mathbb N}.
\]
La correspondencia y su inversa son continuas: cada coordenada finita
depende de un prefijo finito. Los cilindros son abiertos y cerrados.
Cada cilindro de profundidad \(N\) se particiona en sus \(729\)
subcilindros de profundidad \(N+1\); cada uno contiene una prolongación
infinita, obtenida, por ejemplo, continuando con la rama nula en todos
los bloques posteriores.

% Antecedentes efectivos del continuo conjunto y del universo interno.
% Se imprimen en ambos destinos; la etiqueta de universo no reemplaza su regla.
% Punto de inclusión común: después de APP--TRIT--TPK y los bloques 729,
% antes de la cobertura arquimediana y de la carta interna de Feigenbaum.
% No duplica 00a ni incluye clasificación cardinal posterior.
% La regla semántica transfinita se declara como tal, no como consecuencia
% de la sola existencia de prolongaciones finitas.
\begingroup
\let\HMTFGBaseSection\subsection
\let\HMTFGBaseSubsection\subsubsection
\let\HMTFGBaseSubsubsection\paragraph
\let\section\HMTFGBaseSection
\let\subsection\HMTFGBaseSubsection
\let\subsubsection\HMTFGBaseSubsubsection
% Propietario reproducido por unidad demostrativa; véase PROPIETARIOS_CONTINUO_UNIVERSO.json.
% Origen: XI ES CUMPLIMIENTO_EDITORIAL_20260928, sections/continuo_conjunto.tex:1--629.
% Cuerpos y pruebas conservados; sólo namespace, rutas y remisiones contextuales declaradas.
\section{Construcción conjunta del continuo mediante historias y prolongaciones compatibles}
\label{hmtfg:base:iv:continuo-conjunto}
La aritmética expresada por una forma normal conserva una lectura precisa del estado que la produce. Para situar esa lectura dentro de HMT es necesario mantener también las historias, las hojas y las operaciones que actúan sobre ellas. Este capítulo reúne ese sustrato común. Sus distintas construcciones se obtienen del mismo árbol de ejecuciones APP–TRIT–TPK; los cortes de exposición permiten estudiar sus propiedades sin convertirlas en dominios independientes elegidos por una aplicación posterior.

El punto de partida es el núcleo formal expuesto anteriormente. Una emisión admisible contiene la operación de APP, la orientación del TRIT, el cambio de hoja, el residuo y el cociente, junto con la actualización efectiva de memoria del TPK. Cuando se expresa solamente un residuo, la historia de ejecución sigue perteneciendo al objeto de partida. Dos ejecuciones con idéntica lectura final no se identifican por ese solo motivo.

% Propietario reproducido por unidad demostrativa; véase PROPIETARIOS_CONTINUO_UNIVERSO.json.
% Origen: XI ES CUMPLIMIENTO_EDITORIAL_20260928, sections/supervivencia_punto_fijo.tex:1--146.
% Cuerpos y pruebas conservados; sólo namespace, rutas y remisiones contextuales declaradas.
% BEGIN REV09_SUPERVIVENCIA_PUNTO_FIJO
\subsection{Supervivencia por poda de prolongaciones}
\label{hmtfg:base:corpus9:xi:supervivencia}

La memoria retrospectiva distingue las historias registradas por las
operaciones APP--TRIT--TPK; la compatibilidad prospectiva organiza sus
prolongaciones admisibles. Sea \(S\) el espacio de estados completos
alcanzables mediante esas operaciones. Cada estado conserva hojas, residuos,
cocientes, orientación, acarreo, memoria y datos de frontera; cuando la
transición depende de la profundidad, ésta forma parte del estado.
Sea \(E\subseteq S\times S\) la relación efectiva de prolongación y
\(A\subseteq S\) el conjunto admisible determinado por la compatibilidad de
esos registros. Así, la admisibilidad se decide mediante las operaciones y
las condiciones de frontera que generan la historia.

Para \(X\subseteq S\), definimos
\begin{equation}
 \operatorname{Pre}_{E}(X)
 =\{s\in S:\exists t\in X,\ (s,t)\in E\},
 \qquad
 F_{\mathrm{surv}}(X)=A\cap\operatorname{Pre}_{E}(X).
 \label{hmtfg:base:corpus9:xi:operador-poda}
\end{equation}
La sucesión de poda se obtiene por
\begin{equation}
 V_0=A,\qquad
 V_{n+1}=F_{\mathrm{surv}}(V_n),\qquad
 V_\infty=\bigcap_{n\geq0}V_n.
 \label{hmtfg:base:corpus9:xi:sucesion-poda}
\end{equation}

\begin{proposition}[Supervivencia y mayor punto fijo]
\label{hmtfg:base:corpus9:xi:mayor-punto-fijo}
El conjunto \(V_n\) consta exactamente de los estados de \(A\) que admiten
una prolongación de \(n\) pasos cuyos estados pertenecen a \(A\). La
sucesión \((V_n)\) es decreciente. Si cada estado tiene un número finito
de sucesores admisibles, entonces
\[
 V_\infty=F_{\mathrm{surv}}(V_\infty)
          =\nu F_{\mathrm{surv}},
\]
donde \(\nu F_{\mathrm{surv}}\) denota el mayor punto fijo por inclusión.
Sus elementos son precisamente los estados que admiten una prolongación
infinita compatible.
\end{proposition}

\begin{proof}
Para \(n=0\), una prolongación de longitud cero consiste en el propio
estado de \(A\). Si la caracterización vale para \(n\), un estado pertenece
a \(V_{n+1}\) exactamente cuando pertenece a \(A\) y tiene un sucesor en
\(V_n\); anteponer esa arista produce una prolongación admisible de
\(n+1\) pasos. Esta equivalencia prueba la afirmación por inducción.
Además, \(V_1\subseteq V_0\), y la monotonía de
\(F_{\mathrm{surv}}\) propaga \(V_{n+1}\subseteq V_n\).

Sea \(s\in V_\infty\), y escribamos
\(E(s)=\{t\in S:(s,t)\in E\}\). Como \(s\in V_{n+1}\) para todo \(n\),
los conjuntos \(E(s)\cap V_n\) son no vacíos, finitos y decrecientes.
Su intersección contiene un elemento \(t\), que pertenece a
\(V_\infty\) y satisface \((s,t)\in E\). Por tanto,
\(V_\infty\subseteq F_{\mathrm{surv}}(V_\infty)\).
Recíprocamente, la monotonía da
\(F_{\mathrm{surv}}(V_\infty)\subseteq V_{n+1}\) para todo \(n\);
la imagen también está contenida en \(A=V_0\), de modo que pertenece a
\(V_\infty\). Esto prueba la igualdad de punto fijo.

Si \(W=F_{\mathrm{surv}}(W)\), entonces \(W\subseteq A=V_0\).
De \(W\subseteq V_n\) se deduce
\(W=F_{\mathrm{surv}}(W)\subseteq F_{\mathrm{surv}}(V_n)=V_{n+1}\).
La inducción implica \(W\subseteq V_\infty\), que prueba la maximalidad.
Finalmente, el árbol de prolongaciones de un estado de \(V_\infty\)
es finitamente ramificado y tiene un nivel no vacío a toda profundidad.
El lema de König proporciona una rama infinita. En sentido inverso, los
prefijos de cualquier prolongación infinita sitúan su estado inicial
en todos los \(V_n\).
\end{proof}

\paragraph{Ramificación y estabilización.}
La hipótesis de finitud tiene un control explícito. Considérese una raíz
con una rama terminal distinta de cada longitud entera positiva, con todas
las ramas admisibles y disjuntas después de la raíz. La raíz admite
prolongaciones de cualquier longitud finita y pertenece a todos los
\(V_n\); cada sucesor suyo pertenece a una rama de longitud acotada.
Por ello la raíz carece de prolongación infinita y de sucesor en
\(V_\infty\). Este árbol tiene ramificación infinita en la raíz.
En cambio, si \(A\) es finito, cada inclusión estricta
\(V_{n+1}\subsetneq V_n\) elimina al menos un estado. La poda alcanza
un punto fijo después de, como máximo, \(|A|\) eliminaciones estrictas
de capas. Esta cota pertenece al dominio finito \(A\).

\paragraph{Unicidad de la continuación.}
Para \(s\in V_\infty\), los sucesores que conservan supervivencia forman
\[
 E(s)\cap V_\infty.
\]
La selección del próximo estado es única exactamente cuando
\begin{equation}
 \#\bigl(E(s)\cap V_\infty\bigr)=1.
 \label{hmtfg:base:corpus9:xi:unicidad-sucesor}
\end{equation}
Una prolongación infinita es única cuando esta condición se mantiene
en cada estado de la ruta. En efecto, cada sucesor superviviente tiene
una prolongación infinita por la proposición anterior; dos sucesores
distintos producen dos rutas distintas. Cuando el sucesor es único en
cada paso, los prefijos quedan determinados sucesivamente.
La pluralidad de sucesores viables expresa multiplicidad de
prolongaciones dentro del mismo dominio de supervivencia.

\begin{proposition}[Conjugación de las dos direcciones de prolongación]
\label{hmtfg:base:corpus9:xi:conjugacion-supervivencia}
Sea \(C:S\to S\) una involución tal que \(C(A)=A\) y
\begin{equation}
 (x,y)\in E\quad\Longleftrightarrow\quad(Cy,Cx)\in E.
 \label{hmtfg:base:corpus9:xi:involucion-prolongaciones}
\end{equation}
Denotemos por \(V_n^+\) la poda para \(E\) y por \(V_n^-\) la poda
para \(E^{-1}\), ambas con conjunto inicial \(A\). Entonces
\[
 C(V_n^+)=V_n^-\quad(n\geq0),\qquad
 C(V_\infty^+)=V_\infty^-.
\]
\end{proposition}

\begin{proof}
Si \(s_0,\ldots,s_n\) es una ruta admisible para \(E\), la sucesión
\(Cs_0,\ldots,Cs_n\) es una ruta admisible para \(E^{-1}\):
para cada \(j\), la hipótesis convierte
\((s_j,s_{j+1})\in E\) en \((Cs_{j+1},Cs_j)\in E\).
Todos sus estados pertenecen a \(A\). La misma operación aplicada otra
vez recupera la ruta inicial porque \(C^2=I\); por tanto, establece una
bijección entre las prolongaciones finitas de ambas direcciones.
La caracterización de \(V_n^\pm\) por rutas da
\(C(V_n^+)=V_n^-\). Al ser \(C\) biyectiva, conmuta con la
intersección de estos conjuntos, y se obtiene la igualdad de
\(V_\infty^\pm\).
\end{proof}

La involución transporta el conjunto superviviente de una dirección al
de la dirección conjugada. La prolongabilidad infinita en ambas
direcciones corresponde a \(V_\infty^+\cap V_\infty^-\); la unicidad
de una ruta conserva el criterio de
\eqref{hmtfg:base:corpus9:xi:unicidad-sucesor}. La retención de una genealogía
y la restricción de su árbol de continuaciones son dos operaciones
tipadas sobre los mismos registros. La asignación de una entropía a
esas lecturas requiere, además, la medida de historias correspondiente.
% END REV09_SUPERVIVENCIA_PUNTO_FIJO


\subsection{Historias admisibles y memoria por composición}
Sea \(H_0:=\{\ast\}\). El nivel \(H_1\) es el conjunto finito no vacío de
semillas APP--TRIT admisibles, y, para \(k\geq1\), \(H_k\) es el conjunto de
historias enriquecidas de longitud \(k\) que admiten prolongaciones TPK
compatibles a todo horizonte finito. Sus elementos se construyen por emisiones
sucesivas desde el dominio inicial del núcleo. Escribimos \(\zeta\varepsilon\)
para la prolongación de \(\zeta\) por una emisión admisible \(\varepsilon\), y
\(\rho_k(\zeta\varepsilon)=\zeta\) para la supresión de la última emisión. Se
conserva la etiqueta producida por la emisión, incluso cuando diferentes
etiquetas expresan un mismo estado escalar. El espacio de historias completas
queda definido por
\[
 H_\infty^{\rm enr}:=\varprojlim_k(H_k,\rho_k).
\]
El subsistema cofinal \(X_N:=H_{6N}\), con \(X_0=H_0\), coincide con el reloj
de sextetos utilizado en la prueba cardinal; por cofinalidad existe una
identificación canónica
\(H_\infty^{\rm enr}\cong\varprojlim_N(X_N,\rho_{6(N+1),6N})\).

En las dos hojas aritméticas, los datos locales satisfacen las divisiones exactas

\[
\delta_\diamond=r_\diamond+9q_\diamond,
\qquad 0\leq r_\diamond<9,
\qquad \diamond\in\{\Sigma,\Pi\}.
\]
La coordenada orientada del TRIT admite, análogamente, una escritura \(\Delta=o_r+3\kappa\), con el representante orientado declarado por el núcleo. La conservación del cociente permite componer esas representaciones sin sustituir una emisión por su residuo aislado.

En una coordenatización de memoria, la actualización de una emisión se escribe

\[
U_\varepsilon(m)=C_\varepsilon m+\kappa_\varepsilon.
\]
Para una trayectoria \(\alpha=(\varepsilon_1,\ldots,\varepsilon_k)\), la operación efectiva es la composición ordenada de esas actualizaciones. Su parte lineal y su traslación son

\[
C_\alpha=C_{\varepsilon_k}\cdots C_{\varepsilon_1},
\qquad
\kappa_\alpha=
\sum_{j=1}^{k}C_{\varepsilon_k}\cdots
C_{\varepsilon_{j+1}}\kappa_{\varepsilon_j}.
\]
El producto vacío de la última suma es la identidad.

\begin{proposition}[Compatibilidad afín de la concatenación]
\label{hmtfg:base:prop:continuo-compatibilidad-afin}
Si \(\alpha\) precede a \(\beta\), entonces

\[
C_{\beta\alpha}=C_\beta C_\alpha,
\qquad
\kappa_{\beta\alpha}=C_\beta\kappa_\alpha+\kappa_\beta,
\qquad
U_{\beta\alpha}=U_\beta U_\alpha.
\]
\end{proposition}
\begin{proof}
Sustituir \(U_\alpha(m)\) en \(U_\beta\) produce \(C_\beta C_\alpha m+C_\beta\kappa_\alpha+\kappa_\beta\). La descomposición en parte lineal y traslación da las tres identidades. La asociatividad de la composición demuestra que el resultado no depende de una subdivisión auxiliar de la trayectoria.
\end{proof}

Esta identidad de cociclo es el contenido preciso de la memoria acumulada. La cantidad transportada por una segunda trayectoria depende de la actualización de coordenatización que ésta realiza; no es, en general, la suma sin transporte de dos registros.

El retorno nonádico conserva una segunda distinción. En las coordenadas de fase y vuelta, con \(\bar r\in\{0,\ldots,8\}\), su representación es

\[
\sigma_9(w,r)=
\left(w+\left\lfloor\frac{\bar r+1}{9}\right\rfloor,
r+1\pmod9\right),
\qquad
\sigma_9^9(w,r)=(w+1,r).
\]
Durante nueve pasos se produce exactamente un cruce de \(8\) a \(0\). La fase retorna y la coordenada de vuelta aumenta. La holonomía enriquecida retiene además las actualizaciones de memoria que acompañan ese recorrido.

\subsection{Refinamiento del árbol y medida de historias}
La medida inicial es la canónica del censo: \(\mu(\ast)=1\) y
\(\mu(c)=1/|H_1|\) para cada semilla \(c\in H_1\). En este apartado cada
historia tiene un número finito y positivo de prolongaciones. La supresión de
emisiones define las aplicaciones de refinamiento. La admisibilidad del núcleo
proporciona las prolongaciones; cuando una residencia utiliza una emisión
neutra, ésta da explícitamente una sección de \(\rho_k\). La existencia de esa
sección se refiere al árbol que incluye dicha emisión, no a cualquier
restricción arbitraria de rutas.

Sea \(d(\zeta)\geq1\) el número de emisiones admisibles en \(\zeta\). El orden
APP y el censo de prolongaciones determinan la ley canónica
\(p_{\rm can}(\varepsilon\mid\zeta)=1/d(\zeta)\). Sin parámetros
adicionales, esta ley define

\[
\mu(\zeta\varepsilon)=\frac{\mu(\zeta)}{d(\zeta)}.
\]
Una ponderación no uniforme sólo puede sustituir a \(p_{\rm can}\) cuando ya
ha sido producida y añadida al registro enriquecido. El constructor correlativo
de este capítulo utiliza exclusivamente \(p_{\rm can}\); por tanto, su medida
es función del árbol APP--TRIT--TPK y no de una elección externa.

\begin{proposition}[Compatibilidad cilíndrica y esperanza condicional]
\label{hmtfg:base:prop:continuo-esperanza-condicional}
En cada nivel se cumple

\[
\sum_{\rho_k(y)=x}\mu(y)=\mu(x).
\]
En los espacios \(L^2(H_k,\mu_k)\), las aplicaciones

\[
(\mathsf J_kf)(y)=f(\rho_k(y)),
\qquad
(E_kg)(x)=\sum_{\rho_k(y)=x}\frac{\mu(y)}{\mu(x)}g(y)
\]
satisfacen \(E_k=\mathsf J_k^*\), \(E_k\mathsf J_k=I\) y \(\|\mathsf J_kf\|=\|f\|\).
\end{proposition}
\begin{proof}
La primera identidad es la normalización de las probabilidades locales. Aplicándola a \(|f(x)|^2\) se obtiene la isometría. Al agrupar la suma por fibras de \(\rho_k\), se obtiene la fórmula de \(E_k\); sobre una función constante en cada fibra, esa fórmula devuelve su valor.
\end{proof}

Las masas compatibles definen una medida de probabilidad sobre el espacio de historias infinitas: el cilindro de una historia finita recibe su masa \(\mu(\zeta)\). La medida se extiende desde el álgebra de cilindros, y su unicidad procede de que éstos generan la \(\sigma\)-álgebra. Si el árbol es finitamente ramificado y cada historia tiene una prolongación, el límite inverso de sus niveles por supresión de emisiones es no vacío. En la topología cilíndrica, cada cilindro no vacío tiene medida positiva.

Esta medida distingue historias que convergen a una misma representación. La medida uniforme sobre valores terminales sería una lectura distinta, pues eliminaría las multiplicidades que la construcción acaba de conservar.

% Propietario reproducido por unidad demostrativa; véase PROPIETARIOS_CONTINUO_UNIVERSO.json.
% Origen: XI ES CUMPLIMIENTO_EDITORIAL_20260928, sections/ch_memoria_y_naturalidad.tex:1--77.
% Cuerpos y pruebas conservados; sólo namespace, rutas y remisiones contextuales declaradas.
\subsubsection{Lectores enriquecidos y memoria de fibra}
\label{hmtfg:base:sec:ch-memoria-fibra}

La diferencia entre historia y valor admite un criterio operativo. En un
nivel finito \(H_k\), sea \(q_k:H_k\to Z_k\) un lector y sea
\(M_k:H_k\to\mathcal M_k\) el registro conservado. El lector enriquecido
es \(\widehat q_k=(q_k,M_k)\). La memoria distingue estados de una misma
fibra del lector; su suficiencia se comprueba en el dominio considerado.
Esta formulación incorpora al propio artículo el desarrollo completo de la
información de fibra que se utiliza a continuación.

\begin{proposition}[Recuperabilidad en un nivel finito]
\label{hmtfg:base:prop:ch-recuperabilidad-fibra}
Existe una inversa de \(\widehat q_k\) sobre su imagen si y sólo si
\(\widehat q_k\) es inyectivo. Para una variable aleatoria \(Z\) con
valores en \(H_k\) y distribución \(\mu_k\), se cumple
\[
 \mathsf H(Z)=\mathsf H(q_k(Z))+\mathsf H(Z\mid q_k(Z)).
\]
Si \(\widehat q_k\) es inyectivo sobre el soporte de \(Z\), entonces
\[
 \mathsf H(Z\mid q_k(Z),M_k(Z))=0.
\]
Aquí \(\mathsf H(Z)=-\sum_z\Pr(Z=z)\log\Pr(Z=z)\), con
\(0\log0=0\), es la entropía finita en nats; la entropía condicional es
el promedio de la entropía de las distribuciones condicionadas.
\end{proposition}

\begin{proof}
Una inversa recupera un único antecedente, lo que exige inyectividad;
recíprocamente, ésta permite asignar a cada par expresado su único
antecedente. Puesto que \(q_k(Z)\) es función determinista de \(Z\),
agrupar la suma que define \(\mathsf H(Z)\) por fibras de \(q_k\) da la
regla de la cadena. Si el par \((q_k(Z),M_k(Z))\) distingue el soporte,
cada distribución condicionada correspondiente está concentrada en un solo
estado y su entropía es cero.
\end{proof}

No expresar una coordenada y borrar el registro que permite recuperarla son
operaciones distintas. Este resultado no identifica entropía con cardinalidad
del continuo ni introduce una interpretación térmica: su dominio, lector y
distribución son los que se acaban de fijar.

\subsubsection{Equivarianza de la medida bajo simetrías del árbol}
\label{hmtfg:base:sec:ch-naturalidad-medida}

\begin{proposition}[Transporte conjunto del árbol y de su medida]
\label{hmtfg:base:prop:ch-naturalidad-medida}
Sea \(g=(g_k)_{k\geq0}\) una colección de biyecciones \(g_k:H_k\to H_k\)
que conmuta con los truncamientos, fija la raíz y transporta biyectivamente
las emisiones admisibles, conservando sus multiplicidades. Entonces
\[
 d(g_k\zeta)=d(\zeta),\qquad
 (g_k)_*\mu_k=\mu_k,\qquad (g_\infty)_*\mu=\mu.
\]
Más generalmente, si \(\mu_\zeta\) es la ley de continuaciones de una
historia \(\zeta\in H_k\), se tiene
\[
 (g_\infty)_*\mu_\zeta=\mu_{g_k\zeta}.
\]
\end{proposition}

\begin{proof}
La biyección de emisiones conserva su número con multiplicidades. La
distribución uniforme de semillas también es invariante. Cada trayectoria
y su imagen reciben, por tanto, el mismo producto de pesos locales
\(1/d(\zeta)\). Se obtiene la igualdad de medidas en cada nivel y en
cada cilindro. La coherencia con los truncamientos define \(g_\infty\),
y la unicidad de la medida determinada por cilindros prueba la igualdad
en el límite. El mismo argumento, comenzando en \(\zeta\), prueba la
afirmación sobre continuaciones.
\end{proof}

Se obtiene así la naturalidad del árbol y de la medida: un cambio coherente
de representación transporta los pesos y no obliga a escogerlos de nuevo.
La conservación de emisiones y multiplicidades es una hipótesis esencial de
la proposición, no un dato añadido después de la construcción.


\subsection{Transporte por hojas y contabilidad exacta de las salidas}
Consideremos primero el espacio con base ortonormal formada por las historias del nivel \(k\). El transporte que conserva la etiqueta de cada prolongación es

\[
\mathcal U_ke_\zeta=
\sum_{\varepsilon\ \mathrm{admisible\ en}\ \zeta}
\sqrt{p_{\rm can}(\varepsilon\mid\zeta)}\,e_{\zeta\varepsilon}.
\]
Las fibras sucesoras de antecedentes distintos son disjuntas. La normalización de \(p_{\rm can}\) demuestra, por tanto, \(\mathcal U_k^*\mathcal U_k=I\). La representación mediante masas de historias se relaciona con ésta por el cambio de base diagonal que incorpora \(\sqrt{\mu(\zeta)}\).

La etiqueta de hoja forma parte del registro completo de una historia. Se escribe
\[
\widehat\zeta=((\zeta_0,s_0),\ldots,(\zeta_k,s_k)),\qquad
s_j\in\{\Sigma,\Pi\},
\]
donde \(s_j\) es la hoja activa en el paso \(j\). La prolongación añade la emisión y la hoja \(s_{k+1}\) que determina el selector operativo del TPK, conservando íntegro el prefijo, incluida la hoja inicial. En la base de esas historias levantadas,
\[
\mathcal U_ke_{\widehat\zeta}
=\sum_{\varepsilon\ \mathrm{admisible\ en}\ \widehat\zeta}
\sqrt{p_{\rm can}(\varepsilon\mid\widehat\zeta)}\,e_{\widehat\zeta\varepsilon},
\qquad
P_ke_{\widehat\zeta}=\mathbf1_{s_k=\Sigma}e_{\widehat\zeta},\quad
Q_k=I-P_k.
\]
Así, \(P_k\) lee la hoja actual, mientras la distinción entre historias se mantiene por su prefijo completo. Dos historias iniciales que lleguen a la misma hoja activa conservan fibras sucesoras distintas. La prueba de isometría por colecciones disjuntas de prolongaciones se aplica literalmente a esta base.
La orientación trítica es otra coordenada del registro: el valor \(0\) conserva la lectura de umbral en cada hoja y no añade una tercera hoja a esta descomposición. Separamos ahora la conservación de hoja y el intercambio de hoja:

\[
A_k=P_{k+1}\mathcal U_kP_k+Q_{k+1}\mathcal U_kQ_k,
\qquad
\eta_k=Q_{k+1}\mathcal U_kP_k+P_{k+1}\mathcal U_kQ_k.
\]
\begin{proposition}[Identidad de Gram por hojas]
\label{hmtfg:base:prop:continuo-gram-hojas}
Para un transporte lineal \(U\), sea o no isométrico, la descomposición anterior satisface

\[
A^*A+\eta^*\eta=PU^*UP+QU^*UQ.
\]
En particular, para \(U=\mathcal U_k\),

\[
A_k^*A_k+\eta_k^*\eta_k=I.
\]
\end{proposition}
\begin{proof}
En \(A^*A\) y \(\eta^*\eta\), los términos mixtos que contienen \(P_{k+1}Q_{k+1}\) se anulan. Al sumar los términos restantes aparece \(P_kU^*(P_{k+1}+Q_{k+1})UP_k\), y su análogo con \(Q_k\). La primera identidad resulta de la complementariedad de los proyectores; la segunda utiliza además \(U^*U=I\).
\end{proof}

La primera identidad es la fórmula general. Sustituir su lado derecho por \(U^*U\) requiere que los bloques cruzados de ese Gram se anulen. La conservación isométrica del transporte de historias proporciona aquí esa condición; no debe atribuirse automáticamente a una compresión posterior.

Escribamos \(F_{0,0}=I\), \(F_{0,n}=A_{n-1}\cdots A_0\) y \(T_N=F_{0,N}^*F_{0,N}\).

\begin{theorem}[Descomposición finita con residuo terminal]
\label{hmtfg:base:thm:continuo-descomposicion-terminal}
Para cada horizonte \(N\),

\[
I=\sum_{n=0}^{N-1}
F_{0,n}^*\eta_n^*\eta_nF_{0,n}+T_N.
\]
La sucesión \(T_N\) es positiva y decreciente, y posee un límite fuerte positivo \(T_\infty\). La suma de salidas converge fuertemente a \(I-T_\infty\).
\end{theorem}
\begin{proof}
La proposición~\ref{hmtfg:base:prop:continuo-gram-hojas} da \(T_n-T_{n+1}=F_{0,n}^*\eta_n^*\eta_nF_{0,n}\). La suma telescópica prueba la igualdad finita y la monotonía. Para cada vector, las formas cuadráticas decrecen a un límite; la polarización define una forma acotada positiva y, por el teorema de representación de formas acotadas, un operador \(T_\infty\). La convergencia es fuerte porque, para un operador positivo \(S\leq I\), \(\|Sv\|^2\leq\langle v,Sv\rangle\); se aplica a \(S=T_N-T_\infty\).
\end{proof}

El residuo terminal se conserva hasta demostrar su extinción en el dominio considerado. Esta distinción impide que un retorno de fase, por sí solo, se convierta en una hipótesis de agotamiento de todas las historias.

\subsection{Balance orientado y cancelación en el refinamiento}
Una historia lleva su orientación acumulada \(o(h)\) y los recuentos de sus visitas a las dos hojas. El balance local correspondiente es

\[
b(h)=o(h)\bigl(N_\Sigma(h)-N_\Pi(h)\bigr),
\qquad
b^+=\frac{|b|+b}{2},\quad b^-=\frac{|b|-b}{2}.
\]
Cuando un lector expresa la media sobre las prolongaciones, se define \(b_c=Eb_f\). La orientación sigue estando en el argumento de \(b_f\); no se elimina antes de promediar.

\begin{proposition}[Defecto exacto de cancelación]
\label{hmtfg:base:prop:continuo-defecto-cancelacion}
La cantidad

\[
\kappa_{\mathrm{can}}=
\frac{E|b_f|-|Eb_f|}{2}
\]
es no negativa y satisface

\[
E(b_f^+)=(b_c)^++\kappa_{\mathrm{can}},
\qquad
E(b_f^-)=(b_c)^-+\kappa_{\mathrm{can}}.
\]
\end{proposition}
\begin{proof}
La desigualdad triangular ponderada da \(|Eb_f|\leq E|b_f|\). Al escribir las partes positiva y negativa mediante \((|b|\pm b)/2\), se obtienen ambas igualdades.
\end{proof}

La información omitida al expresar solamente el balance es la cancelación común de las contribuciones opuestas. Esta cantidad se determina desde las mismas prolongaciones utilizadas por el transporte anterior. El factor \(1/9\) aparece sólo cuando la fibra contiene nueve prolongaciones con pesos iguales; la esperanza condicional es la regla que sigue siendo válida para ramificación y ponderaciones generales.

\subsection{Incidencia ternaria de las hojas y complejo rectangular}
Las dos hojas determinan el módulo \(V=\mathbb F_3e_\Sigma\oplus\mathbb F_3e_\Pi\). Sus cuatro direcciones proyectivas son las rectas generadas por \(e_\Sigma,e_\Pi,e_\Sigma+e_\Pi,e_\Sigma-e_\Pi\). De ellas proceden seis pares no ordenados. Los ocho vectores no nulos son sus representantes orientados. Así, los recuentos \(4,6,8\) pertenecen a operaciones distintas sobre un mismo módulo de hojas.

Sobre las seis posiciones de pares, definimos

\[
(Ba)_{\{L,L'\}}=a_L+a_{L'},
\qquad
s(q)=\sum_{\{L,L'\}}q_{\{L,L'\}}.
\]
Cada dirección interviene en tres pares; por ello \(sB=0\) en \(\mathbb F_3\). El lector \(W(q)=\mathbf1_{s(q)=0}-1/3\) es invariante bajo \(q\mapsto q+Ba\). Sobre el ambiente uniforme \(\mathbb F_3^6\), sus momentos son

\[
\mathbb EW=0,\qquad
\mathbb EW^2=\frac29,\qquad
\mathbb EW^3=\frac2{27}.
\]
En efecto, \(s\) es sobreyectiva y cada fibra contiene \(3^5\) elementos. El lector vale \(2/3\) con probabilidad \(1/3\) y \(-1/3\) con probabilidad \(2/3\). Estas medias describen el ambiente uniforme; la distribución inducida por un subconjunto de historias debe calcularse con sus propias multiplicidades, ya conservadas en \(H_k\).

Para fijar los puertos sin perder multiplicidades, a profundidad \(k\) se
conservan las lecturas etiquetadas
\(\lambda_\Sigma(\zeta)=(\zeta,\Sigma)\) y
\(\lambda_\Pi(\zeta)=(\zeta,\Pi)\). Sus conjuntos son
\[
\mathscr P_k^\Sigma=\{\lambda_\Sigma(\zeta):\zeta\in H_k\},\qquad
\mathscr P_k^\Pi=\{\lambda_\Pi(\zeta):\zeta\in H_k\}.
\]
Son finitos y no vacíos porque lo es \(H_k\). Los valores, cocientes y residuos
se expresan sobre esos puertos; una igualdad de valores no identifica sus
etiquetas. El producto \(\mathscr P_k^\Sigma\times\mathscr P_k^\Pi\) es el
dominio de la lectura rectangular, dentro del cual se registra el soporte de
las incidencias admisibles.

Fijado este nivel, escribimos \(I=\mathscr P_k^\Sigma\),
\(J=\mathscr P_k^\Pi\) y \(A_m=\mathbb Z/3^m\mathbb Z\). El orden
lexicográfico APP determina la rama superviviente mínima y, en ella, los
puertos de referencia \(i_0,j_0\). No se añade una elección externa: cambiar
de coordenadas produce un complejo canónicamente isomorfo. Las aplicaciones

\[
\kappa(u,v)_{ij}=u_i-v_j,
\qquad
(\delta w)_{ij}=w_{ij}-w_{ij_0}-w_{i_0j}+w_{i_0j_0}
\]
definen la siguiente secuencia, válida sobre el anillo finito \(A_m\):

\[
0\longrightarrow A_m\longrightarrow
A_m^I\oplus A_m^J\mathrel{\mathop{\longrightarrow}^{\kappa}}
A_m^{I\times J}\mathrel{\mathop{\longrightarrow}^{\delta}}
A_m^{(I\setminus\{i_0\})\times(J\setminus\{j_0\})}
\longrightarrow0.
\]
\begin{proposition}[Exactitud de la incidencia rectangular]
\label{hmtfg:base:prop:continuo-exactitud-rectangular}
La secuencia anterior es exacta, con primera flecha dada por la diagonal constante.
\end{proposition}
\begin{proof}
La igualdad \(u_i-v_j=0\) para cada par obliga a que todas las coordenadas de \(u\) y \(v\) sean una misma constante. La sustitución directa da \(\delta\kappa=0\). Si \(\delta w=0\), tómense \(u_i=w_{ij_0}\) y \(v_j=w_{i_0j_0}-w_{i_0j}\). Entonces \(u_i-v_j=w_{ij}\). Finalmente, cualquier matriz del último término se prolonga con fila y columna de referencia nulas; su imagen por \(\delta\) es la matriz inicial. La prueba sólo utiliza sumas y restas, por lo que no exige dividir por una unidad del anillo.
\end{proof}

Reducir de \(A_{m+1}\) a \(A_m\), manteniendo fijos
\(k,\mathscr P_k^\Sigma,\mathscr P_k^\Pi,i_0,j_0\), conmuta con las dos
aplicaciones. Las fórmulas de reconstrucción son las mismas en cada profundidad
modular. Ésta es la compatibilidad aritmética demostrada en este apartado. La
prolongación de historias \(k\to k+1\) conserva las etiquetas de prefijo; sus
aplicaciones entre conjuntos de puertos son un dato distinto del cambio de
anillo \(m\to m+1\).

Esas aplicaciones se obtienen de la supresión de emisiones:
\[
\rho_k^\Sigma(\zeta\varepsilon,\Sigma)=(\zeta,\Sigma),\qquad
\rho_k^\Pi(\zeta\varepsilon,\Pi)=(\zeta,\Pi).
\]
La rama superviviente mínima del orden APP acabada de construir posee prolongaciones compatibles. Sus dos etiquetas definen puertos de referencia compatibles en todos los niveles. El transporte de coordenadas es la precomposición:
\[
(L_ku)_{i'}=u_{\rho_k^\Sigma(i')},\quad
(L_kv)_{j'}=v_{\rho_k^\Pi(j')},\quad
(L_kw)_{i'j'}=w_{\rho_k^\Sigma(i'),\rho_k^\Pi(j')}.
\]
Para el último módulo de la secuencia, se prolonga primero la matriz por cero en la fila y la columna de referencia, se aplica esta misma fórmula y se restringe al nuevo complemento de los puertos de referencia. Esta regla define \(L_k\) también cuando un puerto nuevo tiene como antecedente el puerto marcado.

\begin{proposition}[Naturalidad del complejo rectangular bajo prolongación]
\label{hmtfg:base:iv:rectangular-natural}
Las aplicaciones anteriores verifican
\(L_k\kappa_k=\kappa_{k+1}L_k\) y
\(L_k\delta_k=\delta_{k+1}L_k\).
Además conmutan con la reducción de coeficientes
\(A_{m+1}\to A_m\) y se componen conforme a la supresión de prefijos.
\end{proposition}
\begin{proof}
La primera identidad se obtiene de
\(u_{\rho_k^\Sigma(i')}-v_{\rho_k^\Pi(j')}\).
En la segunda, cada uno de los cuatro términos que define \(\delta\)
se transporta por la misma precomposición. La compatibilidad de los
puertos marcados identifica sus filas y columnas; cuando un antecedente es
el puerto marcado, los cuatro términos se cancelan y coinciden con
la prolongación por cero. Todas las fórmulas son sumas y restas con
coeficientes enteros, de modo que conmutan con la reducción modular.
La composición de precomposiciones es la precomposición con el
prefijo compuesto, lo que demuestra la última afirmación.
\end{proof}

\subsection{Complejo de memoria y compatibilidad de las incidencias}
La misma coordenada acumulada de memoria \(c_y\in\mathbb Z\), en una residencia \(y\), define un complejo libre entero. Su reducción a \(A_m\) utiliza el mismo registro entero antes de reducirlo; por tanto, las profundidades sucesivas son compatibles. Sus generadores son \(f_y\) en grado \(2\), \(e_y,b_y,r_y\) en grado \(1\), y \(g_y\) en grado \(0\):

\[
\partial f_y=3c_y e_y+b_y,\qquad
\partial e_y=g_y,\qquad
\partial b_y=-3c_y g_y,\qquad
\partial r_y=0.
\]
La identidad \(\partial^2=0\) se verifica sobre \(f_y\) mediante \(3c_yg_y-3c_yg_y=0\), y es inmediata en los otros generadores.

Si una trayectoria \(\alpha:y\to y'\) actualiza la memoria, su transporte en el complejo envía los generadores homónimos a sus nuevos representantes, salvo

\[
\Psi_\alpha(b_y)=b_{y'}+3(c_{y'}-c_y)e_{y'}.
\]
\begin{proposition}[Naturalidad del complejo de memoria]
\label{hmtfg:base:prop:continuo-naturalidad-memoria}
Se cumplen \(\partial\Psi_\alpha=\Psi_\alpha\partial\) y \(\Psi_{\beta\alpha}=\Psi_\beta\Psi_\alpha\). Las aplicaciones se prolongan linealmente sobre el anillo de coeficientes elegido.
\end{proposition}
\begin{proof}
Sobre \(f_y\), la imagen de su borde es \(3c_ye_{y'}+b_{y'}+3(c_{y'}-c_y)e_{y'} =3c_{y'}e_{y'}+b_{y'}\). Sobre \(b_y\), el borde de su imagen es \(-3c_{y'}g_{y'}+3(c_{y'}-c_y)g_{y'}=-3c_yg_{y'}\). Los otros dos casos son inmediatos. En una concatenación, las diferencias de memoria se suman telescópicamente, lo que prueba la segunda identidad.
\end{proof}

La relación entre dos representantes de una misma residencia también puede escribirse como un borde explícito. Para una coordenada entera \(v\), reducida en el mismo anillo cuando corresponda, sean

\[
z^-=r_y,\qquad z^+=r_y+3c_yv e_y+v b_y,
\qquad h_*=-vf_y.
\]
Entonces \(\partial h_*=z^--z^+\). Bajo un transporte que actualiza también \(v\), la corrección \(K_\alpha=(v_{y'}-v_y)f_{y'}\) satisface \(K_{\beta\alpha}=K_\beta+\Psi_\beta K_\alpha\). Los símbolos \(K_\alpha\) de esta homotopía local no designan el sello dodecafásico \(K\): pertenecen a otro dominio y su función está especificada por la identidad anterior. Su función puede comprobarse también en la igualdad

\[
z^+_{y'}-\Psi_\alpha z^+_y=\partial K_\alpha.
\]
Ambos miembros valen \((v_{y'}-v_y)(3c_{y'}e_{y'}+b_{y'})\).

\begin{proposition}[Homología y memoria del complejo local]
\label{hmtfg:base:prop:continuo-homologia-memoria}
Sobre \(\mathbb Z\), o después de reducir a \(A_m\), se tiene

\[
H_0(\Gamma)=H_2(\Gamma)=0,
\qquad H_1(\Gamma)\cong A_m[r_y]
\]
en la realización modular, y \(H_1(\Gamma)\cong\mathbb Z[r_y]\) en la realización entera.
\end{proposition}
\Needspace{9\baselineskip}
\begin{proof}
Defínase la homotopía de grado \(1\) mediante \(h(g_y)=e_y\), \(h(b_y)=f_y\), y cero en los otros generadores. Sean \(p\) la proyección al sumando generado por \(r_y\) e \(\iota\) su inclusión. La comprobación sobre cada generador da

\[
\partial h+h\partial=I-\iota p.
\]
En particular, sobre \(b_y\) los términos \(3c_ye_y\) y \(-3c_ye_y\) se cancelan. El complejo se retrae por homotopía al módulo generado por \(r_y\), concentrado en grado \(1\), lo que demuestra la afirmación.
\end{proof}

El transporte \(\Psi_\alpha\) es una cizalla de determinante \(1\). Así, este complejo conserva explícitamente cómo cambia \(c_y\) a nivel de cadenas, pero su homología local no distingue los valores de esa coordenada. La memoria completa permanece en las etiquetas y en los transportes del estado fuente; no se recupera a partir de \(H_1\) aisladamente. Una torsión numérica dependiente de memoria requiere el ensamblaje y la lectura determinantal especificados, además de este complejo local.

\subsection{Ensamblaje con caminos terminales}
Para un horizonte finito, las residencias y sus prolongaciones forman una categoría de caminos. A cada residencia \(\nu\) se asigna el módulo libre \(W(\nu)=A_m[\operatorname{Path}(\nu,\mathrm{Term})]\), que conserva cada continuación terminal como generador. Una trayectoria inicial actúa en \(W\) por precomposición. El complejo \(\Gamma(\nu)\) del apartado anterior actúa covariantemente mediante \(\Psi\).

El ensamblaje utiliza ambos transportes en el mismo módulo bigraduado:

\[
\mathcal B_{q,r}=
\bigoplus_{\nu_0\to\cdots\to\nu_q}
W(\nu_q)\otimes_{A_m}\Gamma_r(\nu_0).
\]
Se utiliza el complejo de caminos normalizado: se omiten las cadenas degeneradas que contienen identidades. El horizonte finito acota entonces el número de prolongaciones estrictas y el grado horizontal. El borde horizontal suprime una residencia intermedia componiendo las flechas adyacentes. En el primer extremo transporta \(\Gamma\); en el último precompone \(W\). Con los signos alternados se obtiene el borde de caminos \(d_{\mathrm{bar}}\). El diferencial total es

\[
D=d_{\mathrm{bar}}+(-1)^q\partial_\Gamma.
\]
Las identidades de caras cancelan por pares los términos de \(d_{\mathrm{bar}}^2\). La naturalidad de la proposición~\ref{hmtfg:base:prop:continuo-naturalidad-memoria} implica que \(d_{\mathrm{bar}}\) conmuta con \(\partial_\Gamma\); el factor \((-1)^q\) cancela los términos cruzados del cuadrado. Por tanto, \(D^2=0\). La compatibilidad se demuestra sobre generadores de caminos y no se introduce por una etiqueta de pertenencia a un complejo.

La evaluación terminal tampoco requiere elegir una historia particular. Para cada continuación \(h\), se aplica la composición efectiva de sus actualizaciones. El resultado conserva el par formado por \(h\) y su lectura, con su multiplicidad. Si \(k<\ell<N\), las continuaciones se descomponen como una unión disjunta de prefijos hasta \(\ell\) y sus continuaciones hasta \(N\). La proposición~\ref{hmtfg:base:prop:continuo-compatibilidad-afin} demuestra que evaluar directamente o evaluar esas dos partes produce la misma lectura. Así se obtiene el cuadrado de naturalidad del terminal.

Los módulos de caminos, las incidencias, el transporte de memoria y sus diferenciales proceden de las mismas prolongaciones. Su ensamblaje no consiste en añadir a un estado cinco coordenadas que se conservarían por definición: sus relaciones se construyen mediante sumas, bordes y composiciones efectivas, y se verifican sobre sus generadores.

\subsection{Determinantes, límite y lectura aritmética}
En cada residencia, el complejo local libre \(\Gamma\) admite su línea determinante graduada

\[
\operatorname{Det}\Gamma=
\bigotimes_r\left(\bigwedge^{\mathrm{rk}\Gamma_r}\Gamma_r\right)^{(-1)^r}.
\]
El ensamblaje total del apartado anterior tiene, por su parte, los módulos \(\operatorname{Tot}_d\mathcal B=\bigoplus_{q+r=d}\mathcal B_{q,r}\). A horizonte finito y sobre el mismo anillo, son libres de rango finito y sólo un número finito de ellos es no nulo. Su línea es
\[
\operatorname{Det}(\operatorname{Tot}\mathcal B)=
\bigotimes_{q,r}
\left(\bigwedge^{\operatorname{rk}\mathcal B_{q,r}}
\mathcal B_{q,r}\right)^{(-1)^{q+r}}.
\]
Esta fórmula resulta de la descomposición por suma directa en cada grado. Distingue la línea del complejo local y la del ensamblaje con todos los caminos; sus bases se ordenan mediante las etiquetas conservadas. La reducción modular conmuta con estas potencias exteriores porque los módulos y sus bases proceden de los mismos módulos libres enteros a horizonte fijado.
La línea está definida por los módulos y sus grados, tanto sobre \(\mathbb Z\) como después de reducir a \(A_m\); el exponente negativo designa el módulo dual. Una torsión numérica no nula exige, adicionalmente, especificar el complejo sobre un cuerpo, las bases y las identificaciones de homología pertinentes. Para calcular factores de Smith se usa la matriz entera producida antes de la reducción, no un levantamiento arbitrario de una matriz modular. Un bloque cuadrado entero de determinante no nulo es invertible sobre \(\mathbb Q\); su reducción es invertible sobre \(A_m\) solamente cuando ese determinante es una unidad, es decir, cuando no es divisible por \(3\). El refinamiento de memorias e incidencias permite formular y comprobar la estabilidad de los factores invariantes pertinentes; estas condiciones no se deducen de la existencia de la línea determinante.

El paso al límite conserva los proyectivos compatibles de residuos, las historias y sus multiplicidades. Los cuadrados probados en los apartados anteriores permanecen conmutativos porque sus aplicaciones están dadas por las mismas fórmulas en cada nivel. Para una coordenada arquimediana, el paso adicional consiste en expresar intervalos cilíndricos no vacíos, compatibles por inclusión y de diámetro tendente a cero. Sus clausuras compactas anidadas tienen una intersección de un solo punto: la existencia se obtiene por compacidad y la unicidad por la cota del diámetro. Ése es el valor de la lectura límite. Si los intervalos son semiabiertos, el punto puede pertenecer sólo a sus clausuras; la regla de representación de frontera determina entonces su escritura posicional. El valor no sustituye la historia que lo ha producido. La otra representación conserva las incidencias del mismo estado. Los dos destinos comparten su origen sin recibir por ello la misma información.

% Procedencia: continuo_conjunto.tex:434--444; remisiones contextuales a
% publicaciones posteriores, no premisas de las pruebas del continuo ni de Feigenbaum.
% Se conserva la descripción positiva sin exigir etiquetas de un anfitrión concreto.
Las regiones coinductivas, el registro
dodecafásico y la incidencia excepcional
se expresan desde ese sustrato con sus mapas
respectivos. La realización hilbertiana conserva después las etiquetas de
fase y especifica su propia inclusión unital; la dualidad y las pantallas
actúan sobre los codominios allí construidos. Las operaciones conjuntas aquí
reunidas transmiten historias, hojas, orientación y memoria a esas
representaciones. La identificación de un funcional concreto con una forma
analítica externa requiere, a su vez, escribir su aplicación sobre los
generadores y verificar sus productos escalares: las identidades de
conservación del presente capítulo no reemplazan esa composición.

\subsection{Alcance matemático de la construcción conjunta}
Las construcciones de estado, árbol y medida; las operaciones correlativas; y
el ensamblaje terminal con naturalidad y límite se desarrollan íntegramente
en este artículo. En este capítulo sus identidades de
balance, incidencia y memoria se han formulado y demostrado sobre los mismos
generadores. Las cinco construcciones permanecen así como operaciones de un
único árbol. Cada especialización conserva visibles sus hipótesis exactas:
isometría del transporte, ley canónica de pesos, tratamiento del residuo
terminal y no singularidad del bloque determinantal cuando corresponda.

Las ecuaciones anteriores prueban los cociclos de memoria, la cancelación
ponderada, la exactitud rectangular y la compatibilidad del complejo de
memoria en todo horizonte finito, además de su paso compatible al límite.
Una identificación analítica posterior debe componer su operador sobre estos
generadores y verificar las propiedades adicionales de su codominio; no
altera ni sustituye la construcción conjunta aquí demostrada.

\iffalse
% Formulación preliminar conservada fuera del ensamblado; sustituida por la
% versión tipada y verificada que se carga al final de este bloque.
\subsection{Emergencia conjunta de las cinco construcciones}
\label{hmtfg:base:sec:cinco-construcciones-continuo-ch}

Las operaciones anteriores se reúnen ahora sin convertirlas en cinco ramas.
Para una historia superviviente \(h\in X_k\), un horizonte \(N>k\) y una
profundidad modular \(m\geq1\), sea el registro base común
\[
 \mathcal B_{k,N}(h):=
 \bigl(\mathscr T_{k,N}(h),
       (\varepsilon_\alpha,s_\alpha,o_\alpha,
        r_{\Sigma,\alpha},q_{\Sigma,\alpha},
        r_{\Pi,\alpha},q_{\Pi,\alpha},
        \operatorname{Carr}_\alpha,
        \kappa_\alpha,\partial_\alpha,
        \operatorname{Surv}_\alpha)_\alpha;
       \operatorname{src},\operatorname{tgt},\circ,\rho\bigr),
 \tag{C1}
 \label{hmtfg:base:eq:base-comun-registros}
\]
donde \(\mathscr T_{k,N}(h)\) es el subárbol de prolongaciones de \(h\)
hasta el horizonte \(N\). Cada componente del registro es una salida de las
divisiones APP, de la orientación TRIT o de las actualizaciones TPK descritas
en las secciones precedentes; ninguna procede de un problema clásico elegido
después.

Sobre esta misma base se definen los cinco objetos estructurados
\[
\begin{aligned}
 \mathcal F_{\rm sp}(\mathcal B_{k,N})
   &:=(\ell^2(X_j,\mu_j),\mathcal U_j,A_j,\eta_j,T_j)_{k\leq j\leq N},\\
 \mathcal F_{\rm sol}(\mathcal B_{k,N})
   &:=(b,\kappa_{\rm can},C_\alpha,\kappa_\alpha)_{\alpha\subset\mathscr T_{k,N}},\\
 \mathcal F_{\rm gau}(\mathcal B_{k,N})
   &:=(I_j,J_j,B_j,\kappa_j,\delta_j)_{k\leq j\leq N},\\
 \mathcal F_{\rm coh}^{(m)}(\mathcal B_{k,N})
   &:=(\Gamma_y\otimes A_m,\partial,\Psi_\alpha,H_\bullet)_{y,\alpha\subset\mathscr T_{k,N}},\\
 \mathcal F_{\rm det}^{(m)}(\mathcal B_{k,N})
   &:=(\operatorname{Tot}\mathcal B,\operatorname{Det}(\operatorname{Tot}\mathcal B))_{k,N,m}.
\end{aligned}
 \tag{C2}
\]
Los discriminantes \({\rm sp},{\rm sol},{\rm gau},{\rm coh},{\rm det}\)
nombran, respectivamente, las lecturas espectral, solenoidal, de calibre,
cohomológica y determinantal de un único objeto. No son dominios de entrada.
Si \(p_T\) olvida la estructura añadida por el lector \(T\) y devuelve el
registro \(\mathcal B_{k,N}(h)\), escribimos
\(\widetilde{\mathcal F}_T=(\mathcal B_{k,N},\mathcal F_T,p_T)\). La salida
correlativa es el producto fibrado
\[
\begin{aligned}
 \mathcal J_{k,N}^{(m)}(h):={}&
 \widetilde{\mathcal F}_{\rm sp}
 \mathop{\times}_{\mathcal B_{k,N}(h)}
 \widetilde{\mathcal F}_{\rm sol}
 \mathop{\times}_{\mathcal B_{k,N}(h)}
 \widetilde{\mathcal F}_{\rm gau}\\
 &\mathop{\times}_{\mathcal B_{k,N}(h)}
 \widetilde{\mathcal F}_{\rm coh}^{(m)}
 \mathop{\times}_{\mathcal B_{k,N}(h)}
 \widetilde{\mathcal F}_{\rm det}^{(m)}.
\end{aligned}
 \tag{C3}
 \label{hmtfg:base:eq:constructor-correlativo}
\]
La igualdad de las cinco proyecciones sobre \(\mathcal B_{k,N}(h)\) obliga a
conservar el mismo árbol, las mismas aristas y el mismo registro; por eso
\eqref{hmtfg:base:eq:constructor-correlativo} no es un producto cartesiano de cinco
objetos ensamblados retrospectivamente.

La supresión de la última capa del árbol induce
\(\operatorname{Res}_{N+1,N}\). Las proposiciones anteriores dan, sobre los
generadores, los cuadrados
\[
\begin{gathered}
 \operatorname{Res}_{N+1,N}\mathcal F_T(\mathcal B_{k,N+1})
 =\mathcal F_T(\mathcal B_{k,N}),
 \qquad T\in\{{\rm sp},{\rm sol},{\rm gau},{\rm coh}\},\\
 \operatorname{Det}(\operatorname{Tot}\mathcal B_{k,N+1})
 \simeq
 \operatorname{Det}(\operatorname{Tot}\mathcal K_{N+1})
 \otimes
 \operatorname{Det}(\operatorname{Tot}\mathcal B_{k,N}),\\
 (\partial\Psi_\alpha=\Psi_\alpha\partial),\qquad
 (L_k\kappa_k=\kappa_{k+1}L_k),\qquad
 (L_k\delta_k=\delta_{k+1}L_k).
\end{gathered}
 \tag{C4}
 \label{hmtfg:base:eq:naturalidad-D}
\]
El cambio \(A_{m+1}\to A_m\) conmuta asimismo con \(\partial\), \(\Psi\),
\(\kappa\), \(\delta\) y las líneas determinantales. Se obtiene, por tanto,
un único prosistema en \((N,m)\).

Sea \(\mathcal X_\chi\) el límite de historias enriquecidas compatibles. Para
\(\xi\in\mathcal X_\chi\), denotemos por \(\mathfrak L_N(\xi)\) el prefijo
ordenado de sus \(N\) emisiones completas. Las fibras conjuntas de horizonte
son
\[
 \mathfrak J_{N,\bullet}^{(m)}
 :=\coprod_{h\in X_0}
 \{(\mathfrak L,J):\mathfrak L\text{ es una rama de longitud }N,
       \ J=\mathcal J_{0,N}^{(m)}(h)\}.
 \tag{C5}
 \label{hmtfg:base:eq:fibras-conjuntas-horizonte}
\]
La marca \(\mathfrak L\) no retiene una copia tautológica de \(\xi\): está
formada por las emisiones y actualizaciones producidas en cada arista. El
generador conjunto es
\[
\begin{aligned}
 \operatorname{Gen}_{\rm cont}:\mathcal X_\chi
   &\longrightarrow\varprojlim_{N,m}\mathfrak J_{N,\bullet}^{(m)},\\
 \xi&\longmapsto
   (\mathfrak L_N(\xi),\mathcal J_{0,N}^{(m)}(h_0))_{N,m},\\
 \mathcal C_{\rm cont}^{\rm disc}
   &:=\operatorname{im}(\operatorname{Gen}_{\rm cont}).
\end{aligned}
 \tag{C6}
 \label{hmtfg:base:eq:continuo-discreto}
\]

\begin{theorem}[Emergencia correlativa de las cinco construcciones]
\label{hmtfg:base:thm:correlacion-cinco}
Para toda historia superviviente \(h\in X_k\), todo horizonte
\(N\geq k+1\) y toda profundidad \(m\geq1\), el producto fibrado
\eqref{hmtfg:base:eq:constructor-correlativo} existe y es no vacío. Sus cinco
componentes emergen del mismo registro APP--TRIT--TPK, satisfacen los
cuadrados de naturalidad \eqref{hmtfg:base:eq:naturalidad-D}, los cambios de base y el
paso al límite. La imagen \(\mathcal C_{\rm cont}^{\rm disc}\) contiene cinco
aspectos inseparables del mismo continuo, no cinco continuos ni cinco
aplicaciones independientes.
\end{theorem}

\begin{proof}
La prolongabilidad del árbol proporciona al menos una arista en cada vértice
interior. El registro de esa arista determina el transporte isométrico y su
descomposición por hojas, el cociclo orientado, los puertos de incidencia, el
complejo de memoria y su línea determinante. Así se construyen sucesivamente
\(\mathcal F_{\rm sp}\), \(\mathcal F_{\rm sol}\),
\(\mathcal F_{\rm gau}\), \(\mathcal F_{\rm coh}^{(m)}\) y
\(\mathcal F_{\rm det}^{(m)}\) sobre la misma base, lo que prueba la no
vacuidad de \eqref{hmtfg:base:eq:constructor-correlativo}. Ninguna de las cinco
construcciones aparece como coordenada de entrada.

La proposición de compatibilidad afín prueba la composición de memoria; la
identidad de Gram, la exactitud rectangular y la naturalidad del complejo
prueban las cuatro primeras líneas de \eqref{hmtfg:base:eq:naturalidad-D}. La
descomposición de la bicomplejidad por la última capa prueba la identidad
multiplicativa de la línea determinante. Todas estas operaciones se calculan
sobre los mismos generadores y conmutan con la reducción modular. La fibra,
la relación, el número, la orientación, el acarreo, la memoria, la frontera y
la ruta permanecen en \(\mathfrak L_N(\xi)\); la multisección, los terminales
y los lectores lineal y probabilístico permanecen en
\(\mathcal J_{0,N}^{(m)}(h_0)\). Por compatibilidad, el límite produce
\eqref{hmtfg:base:eq:continuo-discreto} sin disgregación.
\end{proof}

El reconocimiento ulterior del único continuo construido mediante sus cinco
realizaciones espectral, solenoidal, de calibre, cohomológica y determinantal
pertenece al lenguaje convencional posterior. No selecciona sus estados ni
interviene en su generación.
\fi

% Propietario reproducido por unidad demostrativa; véase PROPIETARIOS_CONTINUO_UNIVERSO.json.
% Origen: XI ES CUMPLIMIENTO_EDITORIAL_20260928, sections/continuo_cinco_construcciones_tipado.tex:1--327.
% Cuerpos y pruebas conservados; sólo namespace, rutas y remisiones contextuales declaradas.
\subsection{Emergencia conjunta de las cinco construcciones}
\label{hmtfg:base:sec:cinco-construcciones-continuo-ch}

Para \(0\leq k\leq N\), sea
\[
 \rho_{N,k}:=\rho_k\circ\rho_{k+1}\circ\cdots\circ\rho_{N-1},
 \qquad \rho_{k,k}:=\operatorname{id}_{H_k},
\]
y, para una historia superviviente \(\zeta\in H_k\), defínase
\[
 \operatorname{Desc}_{k,N}(\zeta)
 :=\{x\in H_N:\rho_{N,k}(x)=\zeta\}.
\]
El árbol finito \(\mathscr T_{k,N}(\zeta)\) tiene como vértices
\[
 \coprod_{j=k}^{N}\{y\in H_j:\rho_{j,k}(y)=\zeta\}
\]
y como aristas únicamente las prolongaciones TPK
\(a:y\to z\) con \(z\in\rho_j^{-1}(y)\). Por tanto, ni sus vértices ni
sus prolongaciones proceden de un árbol auxiliar.

Para \(x\in\operatorname{Desc}_{k,N}(\zeta)\), la rama marcada
\(\lambda_{\zeta,x}\) es la sucesión
\[
 \zeta=x_k\longrightarrow x_{k+1}\longrightarrow\cdots
 \longrightarrow x_N=x
\]
determinada por los truncamientos de \(x\). El registro base común es
\[
\begin{aligned}
 \mathscr R_{k,N}(\zeta;x):=
 \bigl(&\mathscr T_{k,N}(\zeta),\lambda_{\zeta,x};\\
 &\bigl(\varepsilon_a,s_a,o_a,
   (\delta_{\diamond,a},r_{\diamond,a},q_{\diamond,a})
       _{\diamond\in\{\Sigma,\Pi\}},
   \operatorname{Carr}_a,C_a,t_a,\operatorname{Fr}_a,
   \operatorname{Ruta}_a,\operatorname{Mem}_a,
   \operatorname{Surv}_a\bigr)_a;\\
 &\operatorname{src},\operatorname{tgt},\circ,
   (\rho_j)_{k\leq j<N}\bigr),
\end{aligned}
\tag{C1}
\label{hmtfg:base:eq:base-comun-registros}
\]
donde
\[
 U_a(m_0)=C_am_0+t_a,
 \qquad
 \delta_{\diamond,a}=r_{\diamond,a}+9q_{\diamond,a}.
\]
El símbolo \(m_0\) de esta fórmula es una coordenada de memoria; no es la
profundidad modular \(m\) utilizada más abajo. El espacio base verdadero es
\[
 \mathfrak R_{k,N}(\zeta)
 :=\{\mathscr R_{k,N}(\zeta;x):
       x\in\operatorname{Desc}_{k,N}(\zeta)\}.
\]
La rama marcada no es una entrada añadida: es la sucesión de aristas
efectivamente producidas que termina en \(x\). El resto del registro conserva
simultáneamente las fibras laterales del árbol; por ello la marca no sustituye
la ramificación por una trayectoria aislada.

\Needspace{15\baselineskip}
Póngase \(R_m:=\mathbb Z/3^m\mathbb Z\), y sea
\(V_j(\zeta):=\operatorname{Desc}_{k,j}(\zeta)\). El transporte, los
proyectores de hoja y el residuo terminal se restringen al subárbol mediante
\[
 \mathcal U_j^\zeta:=
 \mathcal U_j\!\upharpoonright_{\ell^2(V_j(\zeta),\mu_j)},
 \qquad P_j^\zeta:=P_j\!\upharpoonright_{V_j(\zeta)},
 \qquad Q_j^\zeta:=I-P_j^\zeta,
\]
\[
 A_j^\zeta:=P_{j+1}^\zeta\mathcal U_j^\zeta P_j^\zeta
             +Q_{j+1}^\zeta\mathcal U_j^\zeta Q_j^\zeta,
 \qquad
 \eta_j^\zeta:=Q_{j+1}^\zeta\mathcal U_j^\zeta P_j^\zeta
             +P_{j+1}^\zeta\mathcal U_j^\zeta Q_j^\zeta.
\]
Para \(k\leq n\leq N\), sean
\[
 F_{k,k}^\zeta:=I,
 \qquad
 F_{k,n}^\zeta:=A_{n-1}^\zeta\cdots A_k^\zeta,
 \qquad
 T_{k,n}^\zeta:=(F_{k,n}^\zeta)^*F_{k,n}^\zeta.
\]
Los puertos propios del mismo subárbol son
\[
 \mathscr P_j^\Sigma(\zeta):=
 \{(y,\Sigma):y\in V_j(\zeta)\},
 \qquad
 \mathscr P_j^\Pi(\zeta):=
 \{(y,\Pi):y\in V_j(\zeta)\};
\]
denotamos por \(B_j^\zeta,\kappa_j^\zeta,\delta_j^\zeta\) las restricciones
de la incidencia y del complejo rectangular a esos puertos, y por
\(L_j^\zeta\) la precomposición inducida por la supresión de prefijos.

Para definir sin abreviaturas el lector de caminos, sea
\[
 W_\ell^\zeta(y):=
 R_m[\operatorname{Path}_{\mathscr T_{k,\ell}(\zeta)}
          (y,\operatorname{Term}_\ell)]
\]
y póngase
\[
 \mathscr D_{q,r;k,\ell}^{(m)}[\mathscr R]
 :=\bigoplus_{y_0\to\cdots\to y_q\ \text{en}\
                  \mathscr T_{k,\ell}(\zeta)}
 W_\ell^\zeta(y_q)\otimes_{R_m}
 \bigl(\Gamma_r(y_0)\otimes_{\mathbb Z}R_m\bigr),
 \qquad
 \operatorname{Tot}_d\mathscr D
 :=\bigoplus_{q+r=d}\mathscr D_{q,r;k,\ell}^{(m)}[\mathscr R].
\]
Las fórmulas ya demostradas producen, sobre cada
\(\mathscr R\in\mathfrak R_{k,N}(\zeta)\), los cinco objetos
\[
\begin{aligned}
 \mathcal F_{\rm sp}[\mathscr R]
 &:=(\ell^2(V_j(\zeta),\mu_j),
       \mathcal U_j^\zeta,A_j^\zeta,\eta_j^\zeta,
       T_{k,j}^\zeta)_{k\leq j\leq N},\\
 \mathcal F_{\rm sol}[\mathscr R]
 &:=(b,\kappa_{\rm can},C_\alpha,\kappa_\alpha)
      _{\alpha\subset\mathscr T_{k,N}(\zeta)},\\
 \mathcal F_{\rm gau}^{(m)}[\mathscr R]
 &:=(\mathscr P_j^\Sigma(\zeta),\mathscr P_j^\Pi(\zeta),
       B_j^\zeta,\kappa_j^\zeta,\delta_j^\zeta;R_m)_{k\leq j\leq N},\\
 \mathcal F_{\rm coh}^{(m)}[\mathscr R]
 &:=(\Gamma_y\otimes R_m,\partial,\Psi_\alpha,H_\bullet)
      _{y,\alpha\subset\mathscr T_{k,N}(\zeta)},\\
 \mathcal F_{\rm det}^{(m)}[\mathscr R]
 &:=(\operatorname{Det}(\operatorname{Tot}
       \mathscr D_{\bullet,\bullet;k,\ell}^{(m)}[\mathscr R]))
       _{k\leq\ell\leq N}.
\end{aligned}
\tag{C2}
\label{hmtfg:base:eq:cinco-lectores-comunes}
\]
Aquí \(\mathscr D_{\bullet,\bullet;k,\ell}^{(m)}[\mathscr R]\) es el
bicomplejo normalizado
de caminos y memoria de las secciones anteriores, restringido al subárbol de
horizonte \(\ell\) y reducido sobre \(R_m\). El último lector conserva la
colección ordenada de líneas determinantales hasta el horizonte \(N\); de este
modo, su truncamiento es una proyección canónica y no presupone una
factorización determinantal inexistente. La ley \(\mu\) es la ley canónica
del censo, y los puertos se fijan mediante la rama lexicográficamente mínima
del orden APP. Por tanto, ninguno de los cinco lectores recibe una ponderación
o un puerto externos.

Para
\(T\in\mathfrak T:=\{{\rm sp},{\rm sol},{\rm gau},{\rm coh},{\rm det}\}\),
defínase el espacio total tipado
\[
 \widetilde{\mathcal F}_{T;k,N}^{(m)}(\zeta)
 :=\coprod_{\mathscr R\in\mathfrak R_{k,N}(\zeta)}
   \{\mathcal F_T^{(m)}[\mathscr R]\}_{\mathscr R},
 \qquad
 p_T:\widetilde{\mathcal F}_{T;k,N}^{(m)}(\zeta)
      \longrightarrow\mathfrak R_{k,N}(\zeta),
\]
donde
\(\mathcal F_{\rm sp}^{(m)}:=\mathcal F_{\rm sp}\) y
\(\mathcal F_{\rm sol}^{(m)}:=\mathcal F_{\rm sol}\); es decir, el
superíndice \(m\) es inerte para esos dos lectores.
La proyección \(p_T\) olvida únicamente la estructura del lector y conserva
su registro fuente. El producto fibrado correcto es
\[
\begin{aligned}
 \mathcal J_{k,N}^{(m)}(\zeta):={}&
 \widetilde{\mathcal F}_{\rm sp;k,N}^{(m)}(\zeta)
 \mathop{\times}_{\mathfrak R_{k,N}(\zeta)}
 \widetilde{\mathcal F}_{\rm sol;k,N}^{(m)}(\zeta)
 \mathop{\times}_{\mathfrak R_{k,N}(\zeta)}
 \widetilde{\mathcal F}_{\rm gau;k,N}^{(m)}(\zeta)\\
 &\mathop{\times}_{\mathfrak R_{k,N}(\zeta)}
 \widetilde{\mathcal F}_{\rm coh;k,N}^{(m)}(\zeta)
 \mathop{\times}_{\mathfrak R_{k,N}(\zeta)}
 \widetilde{\mathcal F}_{\rm det;k,N}^{(m)}(\zeta).
\end{aligned}
\tag{C3}
\label{hmtfg:base:eq:constructor-correlativo}
\]
La igualdad de las cinco proyecciones obliga a conservar el mismo árbol, las
mismas aristas, la misma rama producida y el mismo registro. Para cada
\(\mathscr R\) las cinco construcciones determinan el elemento
\[
 s_{k,N}^{(m)}(\mathscr R)
 :=(\mathcal F_T^{(m)}[\mathscr R])_{T\in\mathfrak T}
 \in\mathcal J_{k,N}^{(m)}(\zeta).
\]
La existencia de este elemento procede de las operaciones de
\eqref{hmtfg:base:eq:cinco-lectores-comunes}, no de la
notación del producto fibrado.

La supresión de la última capa induce
\[
 \operatorname{res}^{\mathfrak R}_{N+1,N}
 \bigl(\mathscr R_{k,N+1}(\zeta;x)\bigr)
 :=\mathscr R_{k,N}(\zeta;\rho_N(x)).
\]
La restricción de cada operador a los generadores retenidos define
\[
 r^m_{N+1,N}:\mathcal J_{k,N+1}^{(m)}(\zeta)
 \longrightarrow\mathcal J_{k,N}^{(m)}(\zeta).
\]
En el lector determinantal, \(r^m_{N+1,N}\) suprime únicamente la última
coordenada de la colección de líneas de
\eqref{hmtfg:base:eq:cinco-lectores-comunes}. La reducción
\(R_{m+1}\twoheadrightarrow R_m\) define, por cambio de base,
\[
 q^N_{m+1,m}:\mathcal J_{k,N}^{(m+1)}(\zeta)
 \longrightarrow\mathcal J_{k,N}^{(m)}(\zeta).
\]
Sobre cada generador se cumplen
\[
\begin{gathered}
 r^m_{N+1,N}q^{N+1}_{m+1,m}
 =q^N_{m+1,m}r^{m+1}_{N+1,N},\\
 \partial\Psi_\alpha=\Psi_\alpha\partial,
 \qquad L_j^\zeta\kappa_j^\zeta
       =\kappa_{j+1}^\zeta L_j^\zeta,
 \qquad L_j^\zeta\delta_j^\zeta
       =\delta_{j+1}^\zeta L_j^\zeta,
\end{gathered}
\tag{C4}
\label{hmtfg:base:eq:naturalidad-D}
\]
y las composiciones de los mapas \(r\) y \(q\) satisfacen las identidades
functoriales. La primera igualdad se verifica porque borrar la última capa y
reducir los coeficientes actúan sobre coordenadas distintas. Las otras tres
son las identidades de naturalidad probadas anteriormente. Se obtiene un único
prosistema en \((N,m)\), no cinco límites yuxtapuestos.

En la raíz se define
\[
 \mathfrak J_{N,\bullet}^{(m)}:=\mathcal J_{0,N}^{(m)}(\ast).
\tag{C5}
\label{hmtfg:base:eq:fibras-conjuntas-horizonte}
\]
La estructura discreta conjunta del continuo es el límite de esos objetos
correlativos:
\[
 \boxed{
 \mathcal C_{\rm cont}^{\rm disc}
 :=\varprojlim_{\substack{N\geq1\\m\geq1}}
   \bigl(\mathfrak J_{N,\bullet}^{(m)},
         r^m_{N+1,N},q^N_{m+1,m}\bigr).}
\tag{C6}
\label{hmtfg:base:eq:continuo-discreto}
\]
Ésta es una definición por el prosistema generado por el árbol completo; no es
la imagen ni el grafo de una aplicación que retenga el estado de entrada como
primera coordenada.

Los registros poseen una marca terminal producida. Sus proyecciones compatibles
inducen
\[
 \Pi:\mathcal C_{\rm cont}^{\rm disc}\longrightarrow H_\infty^{\rm enr}.
\]
Recíprocamente, una historia completa determina, por sus prefijos ya emitidos,
las secciones compatibles \(s_{0,N}^{(m)}\). Así se obtiene, después de haber
definido el continuo,
\[
 \operatorname{Gen}_{\rm cont}:H_\infty^{\rm enr}
 \longrightarrow\mathcal C_{\rm cont}^{\rm disc},
 \qquad
 \Pi\circ\operatorname{Gen}_{\rm cont}
 =\operatorname{id}_{H_\infty^{\rm enr}}.
\tag{C7}
\label{hmtfg:base:eq:seccion-historias-continuo}
\]
Esta sección demuestra recuperabilidad; no se usa para definir
\(\mathcal C_{\rm cont}^{\rm disc}\).

\begin{theorem}[Emergencia correlativa de las cinco construcciones]
\label{hmtfg:base:thm:correlacion-cinco}
Para toda historia superviviente \(\zeta\in H_k\), todo \(N>k\), todo
\(m\geq1\) y todo \(\mathscr R\in\mathfrak R_{k,N}(\zeta)\), las cinco
construcciones de \eqref{hmtfg:base:eq:cinco-lectores-comunes} están definidas sobre el
mismo registro, y \eqref{hmtfg:base:eq:constructor-correlativo} contiene
\(s_{k,N}^{(m)}(\mathscr R)\). Sus truncamientos, reducciones modulares y
transportes satisfacen \eqref{hmtfg:base:eq:naturalidad-D}. Además,
\(\mathcal C_{\rm cont}^{\rm disc}\neq\varnothing\), la aplicación \(\Pi\)
es sobreyectiva y queda escindida por \(\operatorname{Gen}_{\rm cont}\). Las
cinco componentes son aspectos inseparables del mismo continuo, no cinco
dominios ni cinco aplicaciones independientes.
\end{theorem}

\begin{proof}
La supervivencia que define \(H_k\) y la ley de prolongación TPK dan
\(\rho_k^{-1}(\zeta)\neq\varnothing\) para cada vértice. Como las fibras son
finitas, el lema de Kőnig produce \(H_\infty^{\rm enr}\neq\varnothing\).

Fijado un registro \(\mathscr R\), el transporte de historias y hojas produce
\(\mathcal F_{\rm sp}[\mathscr R]\); la composición afín y el defecto exacto
de cancelación producen \(\mathcal F_{\rm sol}[\mathscr R]\); las incidencias
y la secuencia rectangular producen \(\mathcal F_{\rm gau}^{(m)}[\mathscr R]\);
el complejo de memoria y su transporte producen
\(\mathcal F_{\rm coh}^{(m)}[\mathscr R]\); y las líneas del ensamblaje
normalizado de caminos producen \(\mathcal F_{\rm det}^{(m)}[\mathscr R]\).
Cada construcción actúa sobre las mismas aristas generadas, no sobre un
subdominio supuesto ni sobre el ambiente abstracto de \(729\) palabras.

Las compatibilidades afín, de transporte, de incidencia y de memoria prueban
\eqref{hmtfg:base:eq:naturalidad-D}. Truncar la colección determinantal de
\eqref{hmtfg:base:eq:cinco-lectores-comunes} conmuta por
definición con \(r\); su cambio de base conmuta con \(q\). Las secciones
\(s_{0,N}^{(m)}\) son compatibles, de modo que producen
\(\operatorname{Gen}_{\rm cont}\). La identidad
\eqref{hmtfg:base:eq:seccion-historias-continuo} prueba a la vez la no
vacuidad del límite, la sobreyectividad de \(\Pi\) y la recuperabilidad de cada
historia sin introducirla como coordenada de definición.

La fibra, la relación, el número, la orientación, el acarreo, la memoria, la
frontera y la ruta permanecen en el registro común; la multisección, los
terminales y los lectores lineal y probabilístico permanecen en las cinco
construcciones correlativas. Por ello el paso al límite conserva la clausura de
dependencias y no disgrega el continuo.
\end{proof}

El reconocimiento ulterior del único continuo construido mediante lenguaje
espectral, solenoidal, de calibre, cohomológico o determinantal pertenece a la
comparación convencional posterior. No selecciona sus estados ni interviene en
su generación APP--TRIT--TPK.


\endgroup
% Propietario reproducido por unidad demostrativa; véase PROPIETARIOS_CONTINUO_UNIVERSO.json.
% Origen: XI ES CUMPLIMIENTO_EDITORIAL_20260928, sections/ch_cierre_cardinal.tex:26--403.
% Cuerpos y pruebas conservados; sólo namespace, rutas y remisiones contextuales declaradas.
\subsection{Prolongación de estados y extensión semántica transfinita}

Sea

\[
 \Sigma_{\rm HMT}
 =\bigl(\dot R_{\rm APP},\dot R_{\rm TRIT},\dot R_{\rm TPK},
 \dot R_{U},\dot R_{\Gamma_9},\dot R_{\rm Ext},\dot R_\rho,
 \dot R_{\pi},\dot R_{\varphi},\dot R_e,\dot R_{\alpha},
 \dot R_{\rm ar},\dot R_{\rm exc}\bigr)
\tag{33}\label{hmtfg:base:eq:ch-signatura}
\]

la signatura \textbf{relacional} numerable cuyo código efectivo es \(h\). Cada
\(\dot R_S\) es el grafo del operador \(S\) sobre los códigos donde está
definido; no se presupone que una función externa sea total dentro de cada
nivel. Los grafos \(\dot R_{\pi},\dot R_{\varphi},\dot R_e,\dot R_{\alpha}\)
son los lectores de salida construidos por la generación coinductiva y el
cierre dodecafásico; no reciben como argumentos los valores convencionales de
las constantes. Su presencia en \(\Sigma_{\rm HMT}\) asegura que la clausura
conserve el estado completo: integra simultáneamente las rutas que expresan
valor y la incidencia que se prolonga hacia la realización excepcional.

En cada horizonte, el código del estado registra las coordenadas APP--TRIT--TPK,
la ruta, la hoja, el acarreo, la frontera y la memoria; en sus etapas de
representación registra asimismo \(P_N,\Phi_N,E_N,A_N\), donde
\(A_N\) es el prefijo de \(\alpha_{\rm HMT}\) producido por el mismo estado.
Estas coordenadas no seleccionan la rama desde fuera: son parte de la historia
que se codifica. Fijada una numeración primitivo-recursiva de las coordenadas finitas de
un estado, toda historia inversa
\(x_\infty=(x_N)_{N\geq1}\in X_\infty^{\rm enr}\) determina el código
\[
 \operatorname{code}(x_\infty)
 :=\{\langle N,j,c\rangle:(x_N)_j\text{ tiene código }c\}\subseteq\omega.
\]
Las ecuaciones de coherencia
\(\rho_{N+1,N}(x_{N+1})=x_N\) permiten decodificar de manera única la
historia a partir de ese real. Escribimos \(m_\infty\) para el código de la
historia realizada y \(m_\infty(n)\) para su función característica bajo la
numeración fijada. El código operatorio \(h\) y este registro completo se
reúnen, mediante una función primitiva recursiva de emparejamiento, en el real

\[
 u_{h,m}:=
 \{\langle0,\langle n,h(n)\rangle\rangle:n<\omega\}
 \cup
 \{\langle1,\langle n,m_\infty(n)\rangle\rangle:n<\omega\}.
\tag{33a}\label{hmtfg:base:eq:ch-codigo-conjunto}
\]

Las dos proyecciones primitivas recuperan \(h\) y \(m_\infty\) desde
\(u_{h,m}\).

La primera extensión es la prolongación métrica y operatoria de los estados
finitos. Para \(N\ge1\), sea
\(\widehat{\mathcal X}^{\rm enr,raw}_{6N}\) el conjunto de estados
enriquecidos brutos de longitud \(6N\), y defínase conjuntistamente el
subsistema superviviente por
\[
\begin{aligned}
 X_N
 &:=\widehat{\mathcal X}^{\rm enr,surv}_{6N}\\
 &:=\Bigl\{x\in\widehat{\mathcal X}^{\rm enr,raw}_{6N}:\\[-.2ex]
 &\qquad\forall M>N\ \exists z\in
   \widehat{\mathcal X}^{\rm enr,raw}_{6M}\;\text{ tal que}\\[-.2ex]
 &\qquad\operatorname{tr}_{6M,6N}(z)=x\ \wedge\
   z\text{ satisface toda transición TPK intermedia}\Bigr\}.
\end{aligned}
\]
y póngase
\[
 \rho_{N+1,N}:=\operatorname{tr}_{6N+6,6N}\!\upharpoonright X_{N+1},
\]
y
\[
 \operatorname{Ext}_N(x):=
 \{y\in X_{N+1}:(x,y)\in
 \operatorname{Ext}_{6N+6,6N}\}.
\]
El nivel raíz se incluye sin excepción implícita:
\[
 X_0:=\{\ast\},\qquad
 \rho_{1,0}:X_1\to X_0\text{ es constante},\qquad
 \operatorname{Ext}_0(\ast):=X_1.
\]
El superíndice \(\mathrm{surv}\) denota el
árbol podado de los estados que poseen descendientes TPK compatibles a todo
horizonte finito. La ley de prolongación TPK establecida en las secciones
precedentes prueba que las semillas
emitidas pertenecen a ese árbol; la ramificación finita y el lema de Kőnig
convierten la compatibilidad a todo horizonte en una rama infinita. En
particular,

\[
 \forall N\ge0\;\forall x\in X_N\quad
 \varnothing\ne\operatorname{Ext}_N(x)
 \subseteq\rho_{N+1,N}^{-1}(x).
\tag{34}\label{hmtfg:base:eq:ch-extension-no-vacia}
\]

Por el mismo lema de árbol finitamente ramificado, toda semilla superviviente
posee una prolongación compatible

\[
 x_\infty=(x_N)_{N\ge1}\in
 X_\infty^{\rm enr}:=\varprojlim_N(X_N,\rho_{N+1,N}).
\tag{35}\label{hmtfg:base:eq:ch-limite-inverso}
\]

La segunda extensión actúa \textbf{después} sobre el código completo de esa
historia realizada.
No se infiere de la mera existencia del operador finito
\(\operatorname{Ext}_N\): formaliza a escala conjuntista el principio HMT de
que sólo son objetos del dominio íntegro los que poseen una derivación
genealógica desde el estado y sus reglas. La extensión semántica transfinita
es:

\[
 \begin{aligned}
  \mathfrak G_0(h,m_\infty)
  &:=\operatorname{TC}\bigl(\{\omega,u_{h,m},h,m_\infty\}\bigr),\\
  \mathfrak G_{\beta+1}(h,m_\infty)
  &:=\operatorname{Def}_{\Sigma_{\rm HMT}}
       \bigl(\mathfrak G_\beta(h,m_\infty)\bigr),\\
  \mathfrak G_\lambda(h,m_\infty)
  &:=\bigcup_{\beta<\lambda}\mathfrak G_\beta(h,m_\infty)
       \quad(\lambda\text{ límite}),\\
  \mathfrak G_{\rm HMT}(h,m_\infty)
  &:=\bigcup_{\beta\in\operatorname{Ord}}
       \mathfrak G_\beta(h,m_\infty).
 \end{aligned}
\tag{36}\label{hmtfg:base:eq:ch-clausura-transfinita}
\]

La hipótesis del continuo y una biyección objetivo no se introducen entre las
condiciones que definen el operador sucesor; el operador está dado
prospectivamente por

\[
 \operatorname{Def}_{\Sigma_{\rm HMT}}(M)
 :=M\cup
 \Bigl\{
   \{x\in M:\langle M,\in,
       (\dot R_S\cap M^{k_S})_{S\in\Sigma_{\rm HMT}}\rangle
       \models\varphi(x,\bar p)\}:
   \ulcorner\varphi\urcorner\in\omega,
   \ \bar p\in M^{<\omega}
 \Bigr\}.
\tag{37}\label{hmtfg:base:eq:ch-operador-def}
\]

Es decir: expresa simultáneamente todos los subconjuntos definibles por una
fórmula finita del lenguaje relacional HMT y parámetros que la genealogía ya
produjo. No
consulta un valor real objetivo, una constante convencional, un cardinal
objetivo ni un modelo exterior. La fórmula~\eqref{hmtfg:base:eq:ch-operador-def} define el cierre transfinito
sobre las semillas, los operadores, el estado enriquecido, el registro y la
doble proyección construidos previamente.

\begin{lemma}[Codificación conjunta del programa HMT y del registro realizado]
\label{hmtfg:base:thm:ch-6}
El par formado por
el código operatorio \(h\) y el registro \(m_\infty\) se codifica por un único
real \(u_{h,m}\subseteq\omega\), y toda fórmula de \(\Sigma_{\rm HMT}\) se
traduce uniformemente a una fórmula del lenguaje \(\{\in,u_{h,m}\}\).

\end{lemma}
\begin{proof}
Tanto \(h\) como \(m_\infty\) son sucesiones de naturales: \(h\)
codifica la signatura, las tablas finitas y los programas de los operadores;
\(m_\infty\) codifica una sola historia realizada, no el espacio completo de
historias. La fórmula~\eqref{hmtfg:base:eq:ch-codigo-conjunto} y sus dos decodificadores prueban la primera
afirmación. Cada símbolo de~\eqref{hmtfg:base:eq:ch-signatura} designa su \textbf{relación de grafo codificada en \(h\)}, nunca el grafo externo de todas sus posibles evaluaciones. APP,
TRIT, \(U_t\), \(\Gamma_9\), las restricciones \(\rho\) y las extensiones
finitas tienen grafos primitivo-recursivos sobre códigos enteros y, por tanto,
fórmulas \(\Delta_0\) absolutas entre niveles transitivos. Las dos
representaciones se expresan como relaciones de grafo mediante las fórmulas

\[
 \begin{aligned}
 \dot R_{\rm ar}(x,y)
 &\;\Longleftrightarrow\;
 \forall k\in\omega\;\exists N\;\forall n\ge N\quad
 |\mu_n(x)-y|<2^{-k},\\
 \dot R_{\rm exc}(x,c,z)
 &\;\Longleftrightarrow\;
 \forall N\in\omega\quad z\!\upharpoonright N
   =E_N(x\!\upharpoonright N,c\!\upharpoonright N),
 \end{aligned}
\tag{37c$_0$}\label{hmtfg:base:eq:ch-grafos-publicacion}
\]

donde \(\mu_n\) y \(E_N\) son las funciones racionales/finitas cuyos códigos
figuran en \(h\). Se sustituye inductivamente cada fórmula atómica que contiene
un símbolo HMT por su fórmula de grafo con parámetro \(u_{h,m}\). Los
conectivos y cuantificadores se conservan. Esta inducción sobre la complejidad
sintáctica produce una traducción \(\varphi\mapsto\varphi^*\) tal que, para
todo nivel transitivo \(M\) de~\eqref{hmtfg:base:eq:ch-clausura-transfinita} que contiene los parámetros,

\[
 \langle M,\in,(\dot R_S\cap M^{k_S})_S\rangle
   \models\varphi(\bar a)
 \quad\Longleftrightarrow\quad
 M\models\varphi^*(\bar a;u_{h,m}).
\tag{37c}\label{hmtfg:base:eq:ch-eliminacion-uniforme}
\]

Así, \eqref{hmtfg:base:eq:ch-operador-def} no puede consultar un oráculo, una tabla decimal objetivo ni una
relación no codificada que introduzca incontables objetos en un solo paso.
\end{proof}

La semántica HMT se define ahora \textbf{independientemente de identificarla con \(\mathfrak G\)}. Sea \(\mathcal T_0^{\rm HMT}\) el conjunto de nombres
canónicos de los elementos de
\(\mathfrak G_0(h,m_\infty)=
\operatorname{TC}(\{\omega,u_{h,m},h,m_\infty\})\). En particular, el nivel
inicial contiene el registro realizado \(m_\infty\), no todas las historias de
\(X_\infty^{\rm enr}\) como parámetros primitivos.
Los términos admisibles son los árboles bien fundados obtenidos por

\[
 \begin{aligned}
 \mathcal T^{\rm HMT}_{\beta+1}
 &:={\mathcal T}^{\rm HMT}_\beta\cup
 \{\langle\beta,\ulcorner\varphi\urcorner,
       t_0,\ldots,t_{k-1}\rangle:
       t_i\in\mathcal T^{\rm HMT}_\beta\},\\
 \mathcal T^{\rm HMT}_\lambda&:=
       \bigcup_{\beta<\lambda}\mathcal T^{\rm HMT}_\beta,
 \qquad
 \mathcal T^{\rm HMT}:=
       \bigcup_{\beta\in\operatorname{Ord}}\mathcal T^{\rm HMT}_\beta.
 \end{aligned}
\tag{37a}\label{hmtfg:base:eq:ch-terminos}
\]

El valor de un término sucesor es el subconjunto del nivel anterior definido
por \(\varphi\) en la estructura relacional inducida, con los valores de
\(t_0,\ldots,t_{k-1}\) como parámetros. Se define, antes de cualquier
comparación conjuntista,

\[
 \operatorname{Ob}_{\rm sem}^{\rm HMT}(h,m_\infty)
 :=\{\llbracket t\rrbracket:t\in\mathcal T^{\rm HMT}\}.
\tag{37b}\label{hmtfg:base:eq:ch-objetos-semanticos}
\]

Ésta es la regla de \textbf{exhaustividad generativa} adoptada por HMT: todo
objeto del dominio íntegro debe poseer uno de esos árboles de derivación y todo
árbol admisible expresa un objeto. No es consecuencia de la ramificación
finita. La hipótesis del continuo no se introduce en esta regla; se obtiene
mediante la demostración desarrollada en el artículo. Tampoco se introducen
como datos una enumeración de los reales ni una biyección objetivo.

\begin{theorem}[Equivalencia entre la semántica de términos y la jerarquía genealógica]
\label{hmtfg:base:thm:ch-6a}
La evaluación de
los términos anteriores satisface

\[
 \operatorname{Ob}_{\rm sem}^{\rm HMT}(h,m_\infty)
 =\mathfrak G_{\rm HMT}(h,m_\infty).
\tag{37d}\label{hmtfg:base:eq:ch-equivalencia-semantica}
\]

\end{theorem}
\begin{proof}
Por inducción transfinita. En el nivel inicial las dos clases
coinciden. Si todo valor de un término de
\(\mathcal T^{\rm HMT}_\beta\) pertenece a
\(\mathfrak G_\beta\), la semántica del término sucesor es uno de los
subconjuntos de~\eqref{hmtfg:base:eq:ch-operador-def}. Recíprocamente, cada subconjunto de~\eqref{hmtfg:base:eq:ch-operador-def} viene dado por
un código de fórmula y una tupla finita de parámetros; por la hipótesis
inductiva esos parámetros poseen términos y~\eqref{hmtfg:base:eq:ch-terminos} forma el término que lo
denota. En un límite ambas construcciones toman la misma unión.
\end{proof}

\Needspace{9\baselineskip}
El codificador de estados es compatible con los truncamientos y envía cada
prolongación TPK a un término sucesor. Sea
\(\mathcal T_N^{\rm st}\subset\mathcal T^{\rm HMT}\) la subclase de términos
que codifica estados completos de profundidad \(N\), y sea
\(\tau_{N+1,N}:\mathcal T_{N+1}^{\rm st}\to\mathcal T_N^{\rm st}\) el borrado
del último bloque junto con el truncamiento correspondiente del registro genealógico,
acarreo y memoria. Si \(J_N:X_N\to\mathcal T_N^{\rm st}\) es el codificador
completo, la compatibilidad de \(\operatorname{Ext}_N\) prueba la identidad
correctamente tipada

\[
 \forall x\in X_N\ \forall y\in\operatorname{Ext}_N(x)\qquad
 \tau_{N+1,N}\bigl(J_{N+1}(y)\bigr)=J_N(x).
\tag{37e$_0$}\label{hmtfg:base:eq:ch-naturalidad-codificador}
\]

Defínase dentro de la semántica
\[
 \operatorname{Succ}^{\rm TPK}_{N+1}(J_N(x)):=
 \{t\in\mathcal T_{N+1}^{\rm st}:\exists y\in X_{N+1}\,
   (y\in\operatorname{Ext}_N(x)\ \wedge\ t=J_{N+1}(y))\}.
\]
Elíjase \(\beta\) tal que \(\mathfrak G_\beta\) contenga \(x\), \(X_N\),
\(X_{N+1}\), los grafos finitos de
\(\operatorname{Ext}_N,J_N,J_{N+1}\) y sus códigos en \(h\). Entonces la
fórmula primitivo-recursiva del grafo de \(\operatorname{Ext}_N\) da

\[
 J_{N+1}[\operatorname{Ext}_N(x)]
 =\operatorname{Succ}^{\rm TPK}_{N+1}(J_N(x)),
 \qquad
 J_{N+1}[\operatorname{Ext}_N(x)]
 \in\mathfrak G_{\beta+1},
\tag{37e}\label{hmtfg:base:eq:ch-sucesores-semanticos}
\]

y ambos miembros son formados por la fórmula de grafo de la misma extensión
TPK. Las ecuaciones~\eqref{hmtfg:base:eq:ch-naturalidad-codificador}--\eqref{hmtfg:base:eq:ch-sucesores-semanticos} son el enlace semántico explícito: la
subclase de términos sucesores es exactamente la imagen de las prolongaciones
admisibles y su truncamiento conmuta. No se identifica la fibra completa de
\(\rho\) con \(\operatorname{Ext}_N\), ni se identifica
\(\operatorname{Def}\) con TPK por nomenclatura.

\begin{lemma}[Transitividad, axiomas y buen orden del modelo genealógico]
\label{hmtfg:base:thm:ch-6b}
\(\mathfrak G_{\rm HMT}(h,m_\infty)\) es una clase transitiva, contiene todos
los ordinales, satisface la teoría de conjuntos de Zermelo--Fraenkel (ZF) y
posee un buen orden global definible; satisface,
por tanto, elección y todos los axiomas usados en la comparación cardinal posterior.

\end{lemma}
\begin{proof}
Abreviemos \(B=\mathfrak G_0\) y \(u=u_{h,m}\). La inducción
sobre~\eqref{hmtfg:base:eq:ch-clausura-transfinita} prueba simultáneamente que cada nivel es transitivo y que
\(\mathfrak G_\alpha\in\mathfrak G_{\alpha+1}\); si todos los ordinales
menores que \(\alpha\) están presentes, su conjunto es definible en el nivel
correspondiente y expresa \(\alpha\). Por ello la unión contiene todos los
ordinales. Extensionalidad y Fundación se heredan del universo ambiente;
Vacío e Infinito están en \(B\). Elegidos un nivel que contenga los parámetros,
las fórmulas que definen \(\{a,b\}\), \(\bigcup a\) y los subconjuntos por
Separación los expresan en un sucesor.

Para Colección--Reemplazo, sea \(F\) una relación funcional definible en
\(\mathfrak G\) sobre el conjunto \(a\). El esquema de reflexión relativo a
la jerarquía definible \((\mathfrak G_\xi,\in,u_{h,m})_{\xi\in\mathrm{Ord}}\),
aplicado en la metateoría ambiente a la lista finita de fórmulas que define
\(F\), produce un nivel
\(\mathfrak G_\theta\) que contiene \(a\), los parámetros y es correcto para
esas fórmulas. Los rangos de los testigos de \(F[a]\) están entonces acotados
por \(\theta\); Separación sobre \(\mathfrak G_\theta\) forma la imagen en
\(\mathfrak G_{\theta+1}\). Para Potencia no se presupone el conjunto interno
que se quiere construir. Fijado \(a\in\mathfrak G_\gamma\), se considera en
la metateoría ambiente el conjunto ya existente \(\mathcal P^V(a)\). Para
cada \(b\in\mathcal P^V(a)\cap\mathfrak G\), sea \(r(b)\) su primera etapa
genealógica. Como \(\mathfrak G\) es una clase definible, Separación en la
metateoría ambiente forma primero el conjunto
\(\mathcal P^V(a)\cap\mathfrak G\). Reemplazo \textbf{en la metateoría ambiente}
aplicado a ese conjunto acota los ordinales \(r(b)\) por algún \(\theta\).
En consecuencia,
\[
 \mathcal P^{\mathfrak G}(a)
 =\{b\in\mathfrak G_\theta:b\subseteq a\},
\]
y el miembro derecho se expresa por Separación en
\(\mathfrak G_{\theta+1}\). No introduce en \(\mathfrak G\) los subconjuntos
exteriores a la semántica~\eqref{hmtfg:base:eq:ch-objetos-semanticos}. Éste es el argumento no circular de
acotación por rangos.

Finalmente, a cada objeto se le asigna su primer nivel de nacimiento y su
menor \textbf{nombre de derivación}: un árbol de derivación bien fundado y
finitamente ramificado (y, por~\eqref{hmtfg:base:eq:ch-terminos}, finito para cada nombre individual), etiquetado por un
ordinal de nivel, un código de fórmula y la tupla de nombres de sus parámetros.
Por recursión transfinita se ordenan esos nombres: primero la etiqueta ordinal,
después el código de fórmula y después, lexicográficamente, los nombres de los
parámetros ya ordenados. El menor nombre de cada valor define una relación de
clase \(<_{\mathfrak G}\) que compara todos los objetos. Es un buen orden
global definible y, tomando el menor elemento de cada conjunto no vacío,
prueba Elección.
\end{proof}

% Propietario reproducido por unidad demostrativa; véase PROPIETARIOS_CONTINUO_UNIVERSO.json.
% Origen: XI ES CUMPLIMIENTO_EDITORIAL_20260928, sections/ch_reflexion_testigos.tex:1--89.
% Cuerpos y pruebas conservados; sólo namespace, rutas y remisiones contextuales declaradas.
\paragraph{Reflexión por clausura de testigos y construcción de la imagen.}
El paso de Colección--Reemplazo del lema~\ref{hmtfg:base:thm:ch-6b} admite la siguiente
demostración explícita. Se mantiene la jerarquía ya generada a partir del
código operatorio y del registro realizado; abreviamos
\(\mathfrak G=\mathfrak G_{\rm HMT}(h,m_\infty)\).
La argumentación utiliza Separación y Reemplazo en la metateoría ambiente,
sin presuponer estos esquemas en \(\mathfrak G\).

Sea \(\Phi\) una colección finita de fórmulas. Se traducen primero sus
cuantificadores universales mediante negación y cuantificación existencial,
y se cierra la colección resultante bajo subfórmulas. Para cualquier nivel
inicial existe un ordinal límite posterior \(\theta\) tal que
\[
 \mathfrak G_\theta\models\varphi(\bar a)
 \quad\Longleftrightarrow\quad
 \mathfrak G\models\varphi(\bar a)
 \qquad
 (\varphi\in\Phi,\ \bar a\in\mathfrak G_\theta^{<\omega}).
\]
La notación de satisfacción de la clase se interpreta, para cada fórmula
fijada, relativizando sus cuantificadores a la clase definible
\(\mathfrak G\); no se introduce un predicado uniforme de verdad.

\begin{proof}
Fijemos \(\alpha\). Para cada subfórmula existencial
\(\exists y\,\psi(y,\bar x)\) de \(\Phi\), las tuplas
\(\bar a\in\mathfrak G_\alpha^{<\omega}\) de la aridad correspondiente
que poseen un testigo en \(\mathfrak G\) forman un conjunto por Separación
metateórica. A cada una se asigna el menor ordinal \(\beta\) para el que
algún \(y\in\mathfrak G_\beta\) satisface
\(\mathfrak G\models\psi(y,\bar a)\). El ordinal existe porque
\(\mathfrak G\) es la unión de sus niveles. Reemplazo en la metateoría
acota esas primeras etapas de testigos. La finitud de \(\Phi\) permite
elegir definiblemente una sola cota \(f(\alpha)>\alpha\), válida para
todas sus subfórmulas existenciales. Por tanto, cada tupla considerada tiene
\emph{al menos un testigo} en \(\mathfrak G_{f(\alpha)}\).

Tomamos \(\alpha_0\) suficientemente grande para contener los parámetros
iniciales y ponemos
\[
 \alpha_{n+1}=f(\alpha_n),\qquad
 \theta=\sup_{n<\omega}\alpha_n.
\]
Por continuidad de la jerarquía, toda tupla finita de
\(\mathfrak G_\theta\) pertenece a algún \(\mathfrak G_{\alpha_n}\).
Si una fórmula existencial de \(\Phi\) con esa tupla es verdadera en
\(\mathfrak G\), posee un testigo en
\(\mathfrak G_{\alpha_{n+1}}\subseteq\mathfrak G_\theta\).
La inducción sobre las fórmulas prueba ahora la equivalencia: las atómicas
se conservan en la estructura inducida; los conectivos preservan la
equivalencia; el paso existencial utiliza ese testigo en una dirección y el
mismo testigo, visto en \(\mathfrak G\), en la otra. La traducción previa
de los universales incluye también las fórmulas existenciales que expresan
sus posibles contraejemplos. Queda probada la reflexión finita requerida.
\end{proof}

Sea ahora \(F(x,y,\bar p)\) funcional sobre \(a\) en \(\mathfrak G\).
Aplicamos la construcción a las fórmulas que expresan \(F\), la existencia
de sus testigos y su imagen, con \(a,\bar p\in\mathfrak G_\theta\).
La transitividad da \(a\subseteq\mathfrak G_\theta\). Cada
\(x\in a\) posee su valor dentro de ese nivel, y la reflexión identifica
la imagen exacta con
\[
 b=\{y\in\mathfrak G_\theta:
       \mathfrak G_\theta\models
       \exists x\in a\ F(x,y,\bar p)\}.
\]
Este conjunto es definible sobre \(\mathfrak G_\theta\), con parámetros
del mismo nivel. Por la operación sucesora de la jerarquía,
\(b\in\operatorname{Def}(\mathfrak G_\theta)
\subseteq\mathfrak G_{\theta+1}\). Así se obtiene Reemplazo en
\(\mathfrak G\). Si se conserva la existencia de testigos y se omite la
unicidad, la misma construcción proporciona un conjunto de testigos para
Colección.

\paragraph{Aridad de los nombres y selección interna.}
En el buen orden del lema~\ref{hmtfg:base:thm:ch-6b}, la enumeración de la base
asigna a cada elemento su menor código natural. En cada sucesor se ordenan
primero los códigos de fórmula; fijado uno, queda fijada la aridad de su
tupla de parámetros. El orden lexicográfico se aplica entonces a un
\emph{producto finito de aridad fija} de los buenos órdenes anteriores.
El menor nombre representa cada valor nuevo y los niveles límite conservan
el orden por primera aparición. De este modo queda especificado el orden
de las tuplas, sin mezclar aridades distintas en un único orden
lexicográfico. La selección del menor elemento de cada conjunto no vacío,
junto con Reemplazo ya obtenido, forma internamente las funciones de
elección utilizadas después. Esta construcción de nombres precede a las
inyecciones cardinales y no recibe ninguna biyección del continuo como
dato.




\subsection{Publicación ordenada de los cilindros de bloque}
\label{hmtfg:publicacion-cilindros-bloque}

Para el representante \([b]_3\in\{0,1,2\}\), se fija
\[
 \nu(\mathbf b)=\sum_{i=1}^6[b_i]_3\,3^{6-i}\in\{0,\ldots,728\}.
\]
En la completación ordenada HMT, un prefijo
\(s=(m;\mathbf b_0,\ldots,\mathbf b_{N-1})\), \(m\in\ZZ\), determina
\[
 a_N(s)=m+\sum_{j=0}^{N-1}\nu(\mathbf b_j)\,729^{-j-1},\qquad
 I_N(s)=[a_N(s),a_N(s)+729^{-N})_{\rm HMT}.
\]
El acarreo conserva el empalme entre los dos representantes de una
frontera. Si \(C_N(s)=\overline{I_N(s)}\), las relaciones anteriores dan
\[
 I_N(s)=\bigsqcup_{\mathbf b\in\FF_3^6}I_{N+1}(s;\mathbf b),
 \qquad \operatorname{diam}C_N(s)=729^{-N}\longrightarrow0.
\]

\begin{proposition}[Cobertura de la lectura arquimediana]
\label{hmtfg:cobertura-arquimediana}
La aplicación
\[
 P_{\rm ar}:\ZZ\times\Omega_\Gamma^{\rm cal}
       \longrightarrow\mathbb R_{\rm HMT},\qquad
 \{P_{\rm ar}(m,\xi)\}
    =\bigcap_N C_N(m;\beta_N(\rho_{\infty,N}\xi))
\]
está bien definida, es sobreyectiva y determina
\[
 (\ZZ\times\Omega_\Gamma^{\rm cal})/\!\sim_{P_{\rm ar}}
       \ \cong\ \mathbb R_{\rm HMT}.
\]
\end{proposition}
\begin{proof}
Los extremos izquierdos forman una sucesión de Cauchy en la completación
HMT, pues los intervalos son anidados y sus diámetros tienden a cero.
Su límite pertenece a todos los intervalos cerrados y es único.
Para cualquier elemento de la completación, se elige primero el intervalo
entero que lo contiene y después el único subintervalo semiabierto de
cada partición que lo contiene. La proposición anterior realiza cada
una de esas elecciones mediante una prolongación TPK efectiva.
Los prefijos compatibles definen \(\xi\), cuya lectura es el elemento
elegido. En las fronteras, el acarreo identifica las escrituras
equivalentes después de su construcción. El cociente por igualdad
de lectura resulta, por definición y sobreyectividad, biyectivo con
la completación.
\end{proof}

La construcción se realiza asimismo en el universo genealógico interno
\(\mathfrak G=\mathfrak G_{\rm HMT}(h,m_\infty)\). Aquí \(h\) codifica
las reglas operatorias y \(m_\infty\) la historia compatible; el universo
se obtiene mediante clausura transitiva inicial, definibilidad en cada
sucesor y unión en los estadios límite:
\[
 G_0=\operatorname{TC}\{\omega,u_{h,m},h,m_\infty\},\qquad
 G_{\beta+1}=\operatorname{Def}_{\Sigma_{\rm HMT}}(G_\beta),\qquad
 G_\lambda=\bigcup_{\beta<\lambda}G_\beta .
\]
Las relaciones de \(\Sigma_{\rm HMT}\) son las de las operaciones ya
construidas y sus códigos. Se mantiene, por tanto, el dominio interno
en todas las afirmaciones de cobertura. Denotemos por
\(\mathbb R_{\rm HMT}^{\mathfrak G}\) su completación y por
\(\mathbb R^{\mathfrak G}\) la realización ordenada dentro del mismo
universo.

\begin{proposition}[Carta ordenada interna]
\label{hmtfg:carta-interna}
Existe un único isomorfismo de cuerpos ordenados arquimedianos completos
\[
 \iota:\mathbb R_{\rm HMT}^{\mathfrak G}
             \xrightarrow{\ \cong\ }\mathbb R^{\mathfrak G}
\]
que fija la unidad. Conserva los extremos racionales, las raíces
cuadradas positivas y los límites de sucesiones convergentes.
\end{proposition}
\begin{proof}
La unidad fija los enteros y los racionales. A cada elemento de la
completación HMT se asocia el corte de los racionales menores que él;
la completitud de la realización ordenada proporciona el único elemento
con ese corte. La densidad racional prueba inyectividad y conservación
estricta del orden. El mismo procedimiento en sentido inverso prueba
sobreyectividad. Las aproximaciones racionales por exceso y por defecto
conservan suma y producto y, con ello, la estructura de cuerpo.
El elemento positivo cuyo cuadrado es \(a>0\) se transporta al elemento
positivo cuyo cuadrado es \(\iota(a)\). Si una sucesión converge, para
cada tolerancia racional positiva sus términos finales están en el
entorno racional del límite; la conservación del orden transporta esa
condición. Todas estas operaciones se efectúan dentro de
\(\mathfrak G\). La unicidad sigue de la conservación de los cortes
racionales.
\end{proof}

\subsection{Descenso por fibras y conservación del selector}
\label{hmtfg:descenso-selector}

% XI REV11, ch_descenso_por_fibras.tex:10--39; prueba completa pertinente.
\begin{proposition}[Criterio de descenso de una operación]
\label{hmtfg:criterio-descenso}
Sea \(q:\Omega\twoheadrightarrow J\) y
\(\star:D\subseteq\Omega^2\to\Omega\), con \(D\) saturado por las
fibras de \(q\times q\). Existe una única operación \(\bar\star\)
sobre \((q\times q)(D)\) tal que
\[
 q(x\star y)=q(x)\,\bar\star\,q(y)
\]
si y sólo si
\[
 q(x)=q(x'),\quad q(y)=q(y')
 \ \Longrightarrow\ q(x\star y)=q(x'\star y')
\]
para los pares del dominio.
\end{proposition}
\begin{proof}
Si existe \(\bar\star\), ambos resultados son la evaluación de la
misma pareja de argumentos. Recíprocamente, se eligen antecedentes
de los argumentos visibles y se define el resultado como la lectura
de su operación. La saturación fija el dominio; la condición sobre
fibras hace independiente el resultado de los antecedentes elegidos.
La sobreyectividad prueba la unicidad.
\end{proof}

En una aplicación mediante \(P_{\rm ar}\), este criterio identifica la
condición precisa para transportar una operación desde estados completos
hasta su lectura. En la carta \(\iota\), la estructura de cuerpo y
los límites ya satisfacen esa condición. A continuación se construye
el perfil analítico a partir del calendario y se utiliza esta
compatibilidad para transportar su iteración y sus raíces.

\subsection{Publicación dodecafásica y normalización cuadrática}
\label{hmtfg:perfil-dodecafasico}

% Propietario integral c37_feigenbaum.tex: comienzo hasta la familia f_lambda.
La sección orientada de criticidad sobre las doce ventanas del calendario
es
\[
 \mathcal P_{\rm crit}(m)=\varphi_m
   =(30m-10)\frac{\pi_{\rm HMT}}{180}
   =\pi_{\rm HMT}\frac{3m-1}{18},\qquad m=1,\ldots,12.
\]
Esta sección fija el origen y la orientación de sus representantes
angulares. El lector uniforme de las doce ventanas es
\[
 \eta_0(F)=\frac1{12}\sum_{m=1}^{12}F(\varphi_m).
\]
Se distingue la función \(\eta_0\), que asigna pesos a estas ventanas,
de la carta ordenada \(\iota\). La normalización uniforme y la sección
angular forman parte de los datos operatorios de este perfil.

Poniendo \(a_m=(3m-1)/18\), se obtiene
\[
 D_0=\frac1{12}\sum_{m=1}^{12}a_m^2=\frac{899}{648},
 \qquad
 s_0=\frac{36}{\pi_{\rm HMT}\sqrt{899}},\qquad
 k_m=s_0\varphi_m=\frac{2(3m-1)}{\sqrt{899}}.
\]
En efecto, \(\sum_{m=1}^{12}(3m-1)^2=5394=6\cdot899\), por lo
que \(s_0^2\pi_{\rm HMT}^2D_0=2\). El perfil es
\begin{equation}
 u_0(x)=\frac1{12}\sum_{m=1}^{12}\bigl[1-\cos(k_mx)\bigr],
 \qquad f_\lambda(x)=1-\lambda u_0(x).
 \label{hmtfg:perfil-familia}
\end{equation}
La constante \(\pi_{\rm HMT}\) entra como publicación angular del
núcleo HMT y se cancela algebraicamente en la normalización del perfil.
La escala \(\lambda\) es el parámetro de la familia, cuyo valor crítico
será seleccionado mediante la dinámica.

\begin{proposition}[Normalización y forma cerrada del perfil]
\label{hmtfg:forma-perfil}
El perfil \(u_0\) es par y analítico, satisface
\(u_0(0)=u_0'(0)=0\) y \(u_0''(0)=2\), y, para
\(y=x/\sqrt{899}\), admite la expresión
\[
 u_0(x)=1-\frac{\cos(37y)\sin(36y)}{12\sin(3y)}
       =1-\frac{\sin(73y)-\sin(y)}{24\sin(3y)}.
\]
Las singularidades aparentes se interpretan por continuidad.
Su expansión inicial es
\[
 u_0(x)=x^2-\frac{715769}{2424603}x^4+O(x^6).
\]
\end{proposition}
\begin{proof}
La suma finita de cosenos conserva paridad y analiticidad. Su derivada
segunda en cero es \(12^{-1}\sum_m k_m^2=2\).
La suma de la progresión trigonométrica
\(\sum_{m=1}^{12}\cos((6m-2)y)\) es
\(\cos(37y)\sin(36y)/\sin(3y)\); la identidad producto--suma
proporciona la segunda expresión. La serie de coseno y
\(\sum_m(3m-1)^4=4294614\) dan el coeficiente de grado cuatro
indicado. La expresión como suma finita fija la extensión en los
ceros del denominador de la forma cerrada.
\end{proof}

\subsection{Transporte de iteradas, raíces y mínimos}
\label{hmtfg:transporte-raices}

% Reunión existente: TRANSVERSALIDAD_Y_ESCALADO_LOCAL_FEIGENBAUM.md, sección 14.1.
Dentro de \(\mathbb R_{\rm HMT}^{\mathfrak G}\), la serie convergente
\[
 \cos_{\rm HMT}(t)=\sum_{j=0}^{\infty}
                   \frac{(-1)^j t^{2j}}{(2j)!}
\]
define la realización analítica correspondiente. Sean
\[
 u_{\rm HMT}(x)=\frac1{12}\sum_{m=1}^{12}
  \left[1-\cos_{\rm HMT}
   \left(\frac{2(3m-1)x}{\sqrt{899}_{\rm HMT}}\right)\right],
 \qquad F_a(x)=1-a\,u_{\rm HMT}(x),
\]
\[
 S_n^{\rm HMT}(a)=F_a^{\,2^n}(0),\qquad
 S_n(\lambda)=f_\lambda^{\,2^n}(0).
\]

\begin{proposition}[Conjugación ordenada de la familia y del selector de raíces]
\label{hmtfg:conjugacion-selector}
Para los argumentos internos de las expresiones anteriores,
\[
 \iota\circ F_a=f_{\iota(a)}\circ\iota,\qquad
 \iota(S_n^{\rm HMT}(a))=S_n(\iota(a)).
\]
La carta transporta exactamente las condiciones de retorno, periodo
exacto y orden de los parámetros. Si un conjunto de raíces admisibles
posee mínimo, la imagen de ese mínimo es el mínimo del conjunto imagen.
\end{proposition}
\begin{proof}
La carta fija los racionales, transporta la raíz cuadrada positiva y
conserva los límites de las sumas parciales de coseno. Por ello,
\(\iota(u_{\rm HMT}(x))=u_0(\iota(x))\). La compatibilidad con
suma y producto prueba la primera igualdad; la inducción sobre el
número de iteraciones prueba la segunda. Una raíz se caracteriza por
la igualdad con cero, que se conserva en ambas direcciones. El
periodo exacto \(2^n\) añade las desigualdades
\(F_a^{2^j}(0)\ne0\), \(0\le j<n\), también conservadas por la
inyectividad de \(\iota\). Si se exige que el parámetro supere un
centro precedente, se transporta esa desigualdad por conservación
del orden. Finalmente, \(a\le b\) para todo \(b\) del conjunto
equivale a \(\iota(a)\le\iota(b)\) para todo elemento de su
imagen. La sobreyectividad sobre los conjuntos así definidos
completa la afirmación acerca del mínimo.
\end{proof}

El dominio de este transporte es el continuo interno indicado.
La existencia de cada mínimo y su identificación con una rama dinámica
concreta son resultados de la construcción de criticidad que sigue;
la carta conserva esas propiedades cuando han sido establecidas.

\subsection{Representación espectral finita del lector uniforme}
\label{hmtfg:jacobi-finito}

% Grassmannianos REV07, sections/jacobi_spectral.tex:35--136,
% mecanismo de Hankel, Gram--Schmidt, recurrencia y representación espectral.
% Especialización atómica ya reunida en la sección 14.2 del desarrollo Feigenbaum.
Las frecuencias cuadradas del perfil definen la medida positiva
\[
 s_m=k_m^2=\frac{4(3m-1)^2}{899},\qquad
 \nu_0=\frac1{12}\sum_{m=1}^{12}\delta_{s_m},\qquad
 M_r=\int s^r\,d\nu_0(s).
\]
Sus doce nodos son distintos y positivos. Para \(1\le r\le12\),
la matriz de Hankel \(H_r=(M_{i+j})_{0\le i,j<r}\) es definida
positiva. En efecto, para \(p\) de grado menor que \(r\),
\[
 \mathbf p^{\mathsf T}H_r\mathbf p
   =\frac1{12}\sum_{m=1}^{12}p(s_m)^2>0
\]
salvo que \(p=0\), porque un polinomio de grado a lo sumo once
que se anula en doce nodos distintos es nulo. En grado doce aparece
el polinomio \(\prod_m(s-s_m)\), de norma nula en esta medida.
La representación espectral correspondiente tiene dimensión doce.

Sean \(p_0,\ldots,p_{11}\) los polinomios mónicos ortogonales,
\(h_j=\int p_j^2\,d\nu_0>0\) y
\(\psi_j=p_j/\sqrt{h_j}\). Como \(\nu_0\) es una probabilidad,
\(\psi_0=1\). La multiplicación por \(s\) en \(L^2(\nu_0)\)
es autoadjunta y posee una matriz tridiagonal \(J_{12}\) en esta base.

\begin{proposition}[Recurrencia y entrelazamiento espectral exactos]
\label{hmtfg:jacobi-entrelazamiento}
Con
\[
 a_j=\frac{\int s p_j(s)^2\,d\nu_0(s)}{h_j},\qquad
 \beta_j=\frac{h_j}{h_{j-1}}>0\quad(1\le j\le11),
\]
la matriz \(J_{12}\) tiene diagonal \(a_j\) y subdiagonal
\(\sqrt{\beta_j}\). Si
\[
 Q_{mj}=\frac{\psi_j(s_m)}{\sqrt{12}}
 \quad(1\le m\le12,\ 0\le j\le11),\qquad
 D=\operatorname{diag}(s_1,\ldots,s_{12}),
\]
entonces
\[
 Q^{\mathsf T}Q=I,\qquad DQ=QJ_{12},\qquad
 J_{12}=Q^{\mathsf T}DQ,
\]
y el perfil satisface
\begin{equation}
 u_0(x)=e_0^{\mathsf T}
          \bigl[I-\cos(x\sqrt{J_{12}})\bigr]e_0.
 \label{hmtfg:perfil-espectral}
\end{equation}
\end{proposition}
\begin{proof}
La positividad de los sistemas de Hankel proporciona de forma única
los polinomios mónicos ortogonales. Al desarrollar \(s p_j\) en la
base ortogonal, los coeficientes de \(p_i\), \(i\le j-2\), se
anulan porque
\(\langle sp_j,p_i\rangle=\langle p_j,sp_i\rangle=0\).
El coeficiente de \(p_{j+1}\) es uno por monicidad, el de \(p_j\)
es \(a_j\) y el de \(p_{j-1}\) es \(h_j/h_{j-1}\).
En el último grado, la misma igualdad se entiende en
\(L^2(\nu_0)\), donde el polinomio mónico de grado doce que se
anula en los nodos representa el vector cero. La normalización da
la subdiagonal positiva indicada.

La ortonormalidad equivale a \(Q^{\mathsf T}Q=I\). La igualdad
de la recurrencia en cada nodo da \(DQ=QJ_{12}\).
Como \(Q\) es cuadrada, también \(QQ^{\mathsf T}=I\), y se
obtiene la diagonalización de \(J_{12}\).
Sus autovalores positivos permiten definir su raíz cuadrada
positiva. Para toda función definida en esos doce nodos,
\[
 e_0^{\mathsf T}\Phi(J_{12})e_0
   =\frac1{12}\sum_{m=1}^{12}\Phi(s_m),
\]
pues \(Qe_0=(1,\ldots,1)^{\mathsf T}/\sqrt{12}\).
La elección \(\Phi(s)=1-\cos(x\sqrt{s})\) prueba
\eqref{hmtfg:perfil-espectral}.
\end{proof}

La representación conserva exactamente el funcional uniforme, sus
momentos y su perfil. Sus primeros coeficientes son
\[
 a_0=M_1=2,\qquad
 \beta_1=M_2-M_1^2=\frac{2493348}{808201}.
\]
La genealogía completa permanece en el estado TPK del que procede
el lector; \(J_{12}\) representa su publicación espectral concreta.
La construcción de Gram--Schmidt y la recurrencia de Jacobi se aplican
a esta medida atómica, con su terminación en dimensión doce.

La cadena operatoria utilizada en las secciones siguientes queda así
fijada: las células de APP y la orientación TRIT alimentan la
transición TPK; sus prolongaciones producen los cilindros y su
lectura ordenada; las doce ventanas orientadas y el funcional uniforme
producen \(u_0\); la familia \(f_\lambda\) determina retornos y raíces;
la carta interna conserva el orden de sus selectores. La realización
espectral describe el mismo perfil mediante una matriz autoadjunta
finita y un vector cíclico.
